\chapter{Integral Satu Dimensi dalam Beberapa Variabel} \label{path:chapter}

%%%%%%%%%%%%%%%%%%%%%%%%%%%%%%%%%%%%%%%%%%%%%%%%%%%%%%%%%%%%%%%%%%%%%%%%%%%%%%

\section{Diferensiasi di bawah tanda integral}
\label{sec:diffunderint}

\sectionnotes{kurang dari satu kuliah}

Misalkan $f(x,y)$ fungsi dua variabel, dan definisikan
\begin{equation*}
g(y) \coloneqq \int_a^b f(x,y) \,dx .
\end{equation*}
Jika $f$ kontinu pada persegi panjang kompak $[a,b] \times [c,d]$, maka
\volIref{\propref*{vI-prop:integralcontcont} dari jilid I}{\propref{prop:integralcontcont}}
menyatakan bahwa $g$ kontinu pada $[c,d]$.

Andaikan $f$ dapat diturunkan terhadap $y$.  
Kapan kita dapat \myquote{mendiferensialkan di bawah tanda integral}?
Dengan kata lain, kapan $g$ diferensiabel dan turunannya diberikan oleh
\begin{equation*}
g'(y) \overset{?}{=} \int_a^b \frac{\partial f}{\partial y}(x,y) \,dx .
\end{equation*}
Diferensiasi merupakan operasi limit, sehingga sebenarnya kita menanyakan kapan
kedua operasi limit berupa integrasi dan diferensiasi dapat dipertukarkan.
Kedua operasi itu tidak selalu dapat dipertukarkan, sehingga diperlukan hipotesis
tambahan.  Pertanyaan pertama
yang kita hadapi adalah keterintegralan
$\frac{\partial f}{\partial y}$, tetapi rumus di atas tetap dapat gagal meskipun
$\frac{\partial f}{\partial y}$ terintegralkan sebagai fungsi dari $x$ untuk setiap
$y$ tetap.

Kita buktikan versi sederhana, tetapi mungkin paling berguna, dari hasil semacam ini.

\begin{thm}[\myindex{Aturan integral Leibniz}]
Andaikan $f \colon [a,b] \times [c,d] \to \R$ fungsi kontinu
sedemikian sehingga $\frac{\partial f}{\partial y}$ ada untuk setiap $(x,y) \in [a,b]
\times [c,d]$ dan kontinu.  Definisikan $g \colon [c,d] \to \R$ dengan
\begin{equation*}
g(y) \coloneqq \int_a^b f(x,y) \,dx .
\end{equation*}
Maka $g$ diferensiabel secara kontinu dan
\begin{equation*}
g'(y) = \int_a^b \frac{\partial f}{\partial y}(x,y) \,dx .
\end{equation*}
\end{thm}

Hipotesis tentang $f$ dan $\frac{\partial f}{\partial y}$ dapat
diperlemah; lihat, misalnya, \exerciseref{exercise:strongerleibniz},
tetapi tidak dapat dihilangkan begitu saja.
Gagasan utama dalam pembuktian mengharuskan
$\frac{\partial f}{\partial y}$ ada dan kontinu untuk semua $x \in [a,b]$,
termasuk kedua titik ujungnya; dalam arah $y$ kita hanya memerlukan
interval kecil.  Dalam penerapan, kita sering mengambil $[c,d]$ sebagai
interval kecil di sekitar titik tempat kita perlu melakukan diferensiasi.

\begin{proof}
Tetapkan $y \in [c,d]$ dan misalkan $\epsilon > 0$ diberikan.
Karena $\frac{\partial f}{\partial y}$ kontinu pada $[a,b] \times [c,d]$, turunan parsial tersebut
kontinu seragam.  Secara khusus, terdapat $\delta > 0$ sedemikian sehingga,
apabila $y_1 \in [c,d]$ dengan
$\sabs{y_1-y} < \delta$, untuk setiap $x \in [a,b]$ kita mempunyai
\begin{equation*}
\abs{\frac{\partial f}{\partial y}(x,y_1)-\frac{\partial f}{\partial y}(x,y)} < \epsilon .
\end{equation*}

Andaikan $h$ memenuhi $y+h \in [c,d]$ dan $0 < \sabs{h} < \delta$.
Tetapkan $x$ sejenak,
lalu terapkan teorema nilai rata-rata untuk memperoleh $y_1$ di antara $y$ dan $y+h$ sedemikian sehingga
\begin{equation*}
\frac{f(x,y+h)-f(x,y)}{h}
=
\frac{\partial f}{\partial y}(x,y_1) .
\end{equation*}
Karena $\sabs{y_1-y} \leq \sabs{h} < \delta$,
\begin{equation*}
\abs{
\frac{f(x,y+h)-f(x,y)}{h}
-
\frac{\partial f}{\partial y}(x,y) 
}
=
\abs{
\frac{\partial f}{\partial y}(x,y_1) 
-
\frac{\partial f}{\partial y}(x,y) 
}
< \epsilon .
\end{equation*}
Argumen tersebut berlaku untuk setiap $x \in [a,b]$ (nilai $y_1$ dapat
bergantung pada $x$).  Dengan demikian, sebagai fungsi dari
$x$,
\begin{equation*}
x \mapsto \frac{f(x,y+h)-f(x,y)}{h}
\qquad
\text{konvergen seragam ke}
\qquad
x \mapsto \frac{\partial f}{\partial y}(x,y)
\qquad
\text{ketika } h \to 0 .
\end{equation*}
Kita mendefinisikan konvergensi seragam untuk barisan, tetapi gagasannya
sama.  Anda dapat mengganti $h$ dengan barisan bilangan tak nol
$\{ h_n \}_{n=1}^\infty$
yang konvergen menuju $0$ sedemikian sehingga $y+h_n \in [c,d]$, lalu membiarkan $n \to \infty$.

Tinjau hasil bagi selisih untuk $g$,
\begin{equation*}
\frac{g(y+h)-g(y)}{h}
=
\frac{\int_a^b f(x,y+h) \,dx -
\int_a^b f(x,y) \,dx }{h}
=
\int_a^b \frac{f(x,y+h)-f(x,y)}{h} \,dx .
\end{equation*}
Konvergensi seragam memungkinkan limit dan integral dipertukarkan.
Jadi,
\begin{equation*}
\lim_{h\to 0}
\frac{g(y+h)-g(y)}{h}
= 
\int_a^b 
\lim_{h\to 0}
\frac{f(x,y+h)-f(x,y)}{h} \,dx 
=
\int_a^b 
\frac{\partial f}{\partial y}(x,y) \,dx .
\end{equation*}
Selanjutnya, berdasarkan
\volIref{\propref*{vI-prop:integralcontcont} dari jilid I}{\propref{prop:integralcontcont}} yang disebutkan di atas,
$g'$ kontinu pada $[c,d]$.
\end{proof}

\begin{example}
Misalkan
\begin{equation*}
f(y) = \int_0^1 \sin(x^2-y^2) \,dx .
\end{equation*}
Maka
\begin{equation*}
f'(y) = \int_0^1 -2y\cos(x^2-y^2) \,dx .
\end{equation*}
\end{example}

\begin{example} \label{example:counterexamplediffunder}
Tinjau
\begin{equation*}
\int_0^{1} \frac{x-1}{\ln(x)} \,dx .
\end{equation*}
Integran tersebut
dapat diperluas menjadi fungsi kontinu pada $[0,1]$, sehingga
integralnya ada; lihat \exerciseref{exercise:counterexamplediffunder}.  Kesulitannya adalah menghitung integral tersebut.
Kita memperkenalkan parameter $y$
dan mendefinisikan suatu fungsi:
\begin{equation*}
g(y) \coloneqq \int_0^{1} \frac{x^y-1}{\ln(x)} \,dx .
\end{equation*}
Fungsi
$\frac{x^y-1}{\ln(x)}$
juga dapat diperluas menjadi fungsi kontinu pada
$[0,1] \times [0,1]$ (ini juga bagian dari latihan tersebut).
Lihat \figureref{fig:diffunderexample}.
\begin{myfigureht}
\myincludegraphics{diffunderexample}{%
Plot permukaan suatu fungsi di kuadran pertama pada bidang xy.
Fungsi bernilai nol pada sumbu x maupun sumbu y, tetapi menanjak sangat
curam ketika menjauh dari sumbu y.}
\caption{Grafik $z= \frac{x^y-1}{\ln(x)}$ pada $[0,1] \times [0,1]$.\label{fig:diffunderexample}}
\end{myfigureht}


Oleh karena itu,
$g$ merupakan fungsi kontinu pada $[0,1]$ dan $g(0) = 0$.
Untuk setiap $0 < \epsilon < 1$, turunan parsial integran terhadap $y$, yaitu $x^y$,
memiliki perluasan kontinu pada $[0,1] \times [\epsilon,1]$.  Oleh karena itu,
untuk $0 < y \leq 1$, kita dapat mendiferensialkan di bawah tanda integral,
\begin{equation*}
g'(y) =
\int_0^{1} \frac{\ln(x) x^y}{\ln(x)} \,dx 
=
\int_0^{1} x^y \,dx =
\frac{1}{y+1} .
\end{equation*}
Pada baris pertama, pembatalan faktor $\ln(x)$ berlaku untuk $0<x<1$; nilai-nilai
titik ujung diberikan oleh perluasan kontinu dan tidak mengubah integral Riemann.
Kita perlu menghitung $g(1)$ dengan mengetahui bahwa $g'(y) = \frac{1}{y+1}$ untuk
$0<y\leq1$ dan $g(0) = 0$.  Kalkulus elementer menyatakan bahwa
$g(1) = \int_0^1 g'(y)\,dy = \ln(2)$.
Dengan demikian,
\begin{equation*}
\int_0^{1} \frac{x-1}{\ln(x)} \,dx  = \ln(2).
\end{equation*}
\end{example}

\subsection{Latihan}

\begin{exercise} \label{exercise:counterexamplediffunder}
Buktikan dua pernyataan yang dinyatakan dalam
\exampleref{example:counterexamplediffunder}:
\begin{enumerate}[a)]
\item
Buktikan bahwa $\frac{x-1}{\ln(x)}$ dapat diperluas menjadi fungsi kontinu pada
$[0,1]$.  Artinya, terdapat fungsi kontinu pada $[0,1]$
yang sama dengan $\frac{x-1}{\ln(x)}$ pada $(0,1)$.
\item
Buktikan bahwa $\frac{x^y-1}{\ln(x)}$ dapat diperluas menjadi fungsi kontinu
pada $[0,1] \times [0,1]$.
\end{enumerate}
\end{exercise}

\begin{exercise}
Andaikan $h \colon \R \to \R$ kontinu dan $g
\colon \R \to \R$ diferensiabel secara kontinu serta berdukungan
kompak.  Artinya, terdapat suatu $M > 0$ sedemikian sehingga $g(x) = 0$ apabila
$\sabs{x} \geq M$.  Definisikan
\begin{equation*}
f(x) \coloneqq \int_{-\infty}^\infty h(y)g(x-y)\,dy  .
\end{equation*}
Buktikan bahwa $f$ diferensiabel.
\end{exercise}

\begin{exercise}
Andaikan $f \colon \R \to \R$ diferensiabel tak berhingga kali (turunan
semua orde ada)
dan $f(0) = 0$.  Buktikan bahwa terdapat fungsi $g \colon \R \to \R$ yang
diferensiabel tak berhingga kali sedemikian sehingga $f(x) = x\,g(x)$.
Buktikan juga bahwa
jika $f'(0) \neq 0$, maka $g(0) \neq 0$.\\
Petunjuk: Tuliskan
$f(x) = \int_0^x f'(s) \,ds$, lalu ubah batas integralnya menjadi $0$ dan $1$.
\end{exercise}

\begin{exercise}
Hitung $\int_0^1 e^{tx} \,dx$.  Untuk setiap bilangan bulat $n \geq 0$, peroleh rumus untuk
$\int_0^1 x^n e^{x} \,dx$ tanpa menggunakan integrasi parsial, melainkan
melalui diferensiasi di bawah tanda integral.
\end{exercise}

\begin{exercise}
Misalkan $U \subset \R^n$ terbuka dan andaikan
$f(x,y_1,y_2,\ldots,y_n)$ adalah fungsi kontinu
yang didefinisikan pada $[0,1] \times U \subset \R^{n+1}$.
Andaikan
$\frac{\partial f}{\partial y_1},
\frac{\partial f}{\partial y_2},\ldots,
\frac{\partial f}{\partial y_n}$
ada dan kontinu pada $[0,1] \times U$.
Buktikan bahwa $F \colon U \to \R$ yang didefinisikan oleh
\begin{equation*}
F(y_1,y_2,\ldots,y_n) \coloneqq
\int_0^1
f(x,y_1,y_2,\ldots,y_n)
\, dx
\end{equation*}
diferensiabel secara kontinu.
\end{exercise}

\begin{myfigureht}
\myincludegraphics{diffunderex916}{%
Plot permukaan suatu fungsi pada persegi satuan.  Pada kedua sumbu fungsi
bernilai nol; di bagian lain persegi fungsi bernilai positif.  Namun, di titik asal
grafik tampak tidak kontinu: ketika bergerak dari bagian dalam persegi menuju
titik asal, nilai fungsi tampaknya mendekati suatu nilai positif, bukan nol.}
\caption{Grafik $z= \frac{xy^3}{{(x^2+y^2)}^2}$ pada
$[0,1] \times [0,1]$.\label{fig:diffunderex916}}
\end{myfigureht}

\begin{exercise}
\pagebreak[2]
Kerjakan secara rinci contoh tandingan berikut.  Misalkan
\begin{equation*}
f(x,y) \coloneqq
\begin{cases}
\frac{xy^3}{{(x^2+y^2)}^2} & \text{jika } x \neq 0 \text{ atau } y \neq 0, \\
0                          & \text{jika } x = 0 \text{ dan } y = 0.
\end{cases}
\end{equation*}
Lihat \figureref{fig:diffunderex916}.
\begin{enumerate}[a)]
\item
Buktikan bahwa untuk setiap $y$ tetap, fungsi $x \mapsto f(x,y)$
terintegralkan Riemann pada $[0,1]$, dan
\begin{equation*}
g(y) \coloneqq \int_0^1 f(x,y) \, dx = \frac{y}{2(y^2+1)} .
\end{equation*}
Dengan demikian, $g'(y)$ ada dan merupakan fungsi kontinu:
\begin{equation*}
g'(y) =
\frac{d}{dy} \int_0^1 f(x,y) \, dx
=
\frac{1-y^2}{2{(y^2+1)}^2} .
\end{equation*}
\item
Buktikan bahwa $\frac{\partial f}{\partial y}$ ada untuk semua $x$ dan $y$, lalu
hitunglah turunan tersebut.
\item
Buktikan bahwa untuk semua $y$, integral
\begin{equation*}
\int_0^1 \frac{\partial f}{\partial y} (x,y) \, dx
\end{equation*}
ada, tetapi
\begin{equation*}
g'(0) \neq \int_0^1 \frac{\partial f}{\partial y} (x,0) \, dx .
\end{equation*}
\end{enumerate}
\end{exercise}

\begin{exercise}
\pagebreak[2]
Kerjakan secara rinci contoh tandingan berikut.  Misalkan
\begin{equation*}
f(x,y) \coloneqq
\begin{cases}
x \,\sin \left(\frac{y}{x^2+y^2}\right) & \text{jika } (x,y) \neq (0,0),\\
0                                       & \text{jika } (x,y)=(0,0).
\end{cases}
\end{equation*}
\begin{enumerate}[a)]
\item
Buktikan bahwa $f$ kontinu pada seluruh $\R^2$.
Dengan demikian, fungsi berikut terdefinisi dengan baik untuk setiap $y \in \R$:
\begin{equation*}
g(y) \coloneqq \int_0^1 f(x,y) \, dx .
\end{equation*}
\item
Buktikan bahwa $\frac{\partial f}{\partial y}$ ada untuk semua $(x,y)$,
tetapi tidak kontinu di $(0,0)$.
\item
Buktikan bahwa $\int_0^1 \frac{\partial f}{\partial y}(x,0) \, dx$ tidak
ada bahkan jika kita menggunakan integral tak wajar; artinya,
limit
$\lim\limits_{h \to 0^+} \int_h^1 \frac{\partial f}{\partial y}(x,0) \, dx$
tidak ada.
\end{enumerate}
Catatan: Anda boleh menggunakan pengetahuan tentang sinus dan kosinus dari kalkulus.
\end{exercise}

\begin{exercise} \label{exercise:strongerleibniz}
\pagebreak[3]
Perkuat aturan integral Leibniz sebagai berikut.
Andaikan $f \colon (a,b) \times (c,d) \to \R$ adalah fungsi kontinu yang terbatas
sedemikian sehingga $\frac{\partial f}{\partial y}$ ada untuk semua $(x,y) \in (a,b)
\times (c,d)$ serta kontinu dan terbatas.  Definisikan $g \colon (c,d) \to \R$
dengan
\begin{equation*}
g(y) \coloneqq \int_a^b f(x,y) \,dx .
\end{equation*}
Integral di atas, demikian pula integral pada rumus berikut, dipahami sebagai
integral tak wajar pada kedua titik ujung.
Maka $g$ diferensiabel secara kontinu dan
\begin{equation*}
g'(y) = \int_a^b \frac{\partial f}{\partial y}(x,y) \,dx .
\end{equation*}
Petunjuk: Lihat juga \volIref{\exerciseref*{vI-exercise:integralcontcontextra} dan
\thmref*{vI-thm:dersconverge} dari
jilid I}{\exerciseref{exercise:integralcontcontextra} dan
\thmref{thm:dersconverge}}.
\end{exercise}

\begin{exercise}
Andaikan $g \colon \R \to \R$ diferensiabel secara kontinu,
$h \colon \R^2 \to \R$ kontinu, dan
$\frac{\partial h}{\partial x}$ ada serta kontinu di setiap titik.  Buktikan bahwa
\begin{equation*}
F(x,y) \coloneqq g(x) + \int_0^y h(x,s) \,ds
\end{equation*}
diferensiabel secara kontinu dan merupakan solusi persamaan diferensial parsial
$\frac{\partial F}{\partial y} = h$ dengan syarat awal
$F(x,0) = g(x)$ untuk semua $x \in \R$.
\end{exercise}

%%%%%%%%%%%%%%%%%%%%%%%%%%%%%%%%%%%%%%%%%%%%%%%%%%%%%%%%%%%%%%%%%%%%%%%%%%%%%%

\sectionnewpage
\section{Integral lintasan}
\label{sec:pathintegral}

%mbxINTROSUBSECTION

\sectionnotes{2--3 kuliah}

\subsection{Lintasan mulus sepotong-sepotong}

Misalkan $\gamma \colon [a,b] \to \R^n$ suatu fungsi dan tuliskan
$\gamma = (\gamma_1,\gamma_2,\ldots,\gamma_n)$. 
Andaikan $\gamma$ 
\emph{diferensiabel secara kontinu},
yang berarti bahwa fungsi tersebut diferensiabel dan turunannya kontinu.
Dengan kata lain, terdapat fungsi kontinu $\gamma' \colon [a,b]
\to \R^n$ sedemikian sehingga, untuk setiap $t \in [a,b]$, berlaku
$\lim\limits_{h \to 0}
\frac{\snorm{\gamma(t+h)-\gamma(t) - \gamma'(t) \, h}}{\sabs{h}} = 0$.
Pada titik ujung, limit ini dipahami secara relatif terhadap $[a,b]$,
sehingga $t+h \in [a,b]$.
Kita memandang
$\gamma'(t)$ sebagai operator linear (matriks $n \times 1$) atau sebagai
vektor,
$\gamma'(t) =
\bigl( \gamma_1'(t), \gamma_2'(t), \ldots, \gamma_n'(t) \bigr)$.
Secara ekuivalen, 
$\gamma_j$ adalah fungsi yang diferensiabel secara kontinu pada $[a,b]$
untuk setiap $j=1,2,\ldots,n$.
Menurut \exerciseref{exercise:normonedim}, norma operator dari
operator $\gamma'(t)$ sama dengan
norma Euklides dari vektor yang bersesuaian, sehingga kita dapat menuliskan 
$\bnorm{\gamma'(t)}$ tanpa menimbulkan kerancuan.

\begin{defn}
Fungsi $\gamma \colon [a,b] \to \R^n$
disebut \emph{\myindex{lintasan mulus}}
atau
\emph{\myindex{lintasan yang diferensiabel secara kontinu}}\footnote{Dalam
literatur, kata \myquote{mulus} kadang-kadang berarti
\myquote{diferensiabel tak berhingga kali}.}
jika
$\gamma$ diferensiabel secara kontinu dan
$\gamma'(t) \neq 0$ untuk semua $t \in [a,b]$.

Fungsi $\gamma \colon [a,b] \to \R^n$ disebut
\emph{\myindex{lintasan mulus sepotong-sepotong}} atau
\emph{\myindex{lintasan diferensiabel secara kontinu sepotong-sepotong}}
jika terdapat berhingga banyak titik
$t_0 = a < t_1 < t_2 < \cdots < t_k = b$ sedemikian sehingga
pembatasan $\gamma|_{[t_{j-1},t_j]}$ merupakan lintasan mulus untuk setiap
$j=1,2,\ldots,k$.

Lintasan $\gamma$ disebut 
\emph{\myindex{lintasan tertutup}} jika $\gamma(a) = \gamma(b)$; artinya,
lintasan tersebut bermula dan berakhir di titik yang sama.
Lintasan $\gamma$ disebut \emph{\myindex{lintasan sederhana}} jika 
1) $\gamma$ merupakan fungsi injektif, atau 2)
$\gamma|_{[a,b)}$ injektif dan $\gamma(a)=\gamma(b)$ ($\gamma$ merupakan
lintasan tertutup sederhana).
\end{defn}


\begin{example} \label{mv:example:unitsquarepath}
Misalkan $\gamma \colon [0,4] \to \R^2$ didefinisikan oleh
\begin{equation*}
\gamma(t) \coloneqq
\begin{cases}
(t,0)   & \text{jika } t \in [0,1],\\
(1,t-1) & \text{jika } t \in (1,2],\\
(3-t,1) & \text{jika } t \in (2,3],\\
(0,4-t) & \text{jika } t \in (3,4].
\end{cases}
\end{equation*}
\begin{myfigureht}
\myincludegraphics{squarepath}{%
Diagram persegi pada bidang xy dengan panah-panah yang menunjukkan arah lintasan.
Pada t sama dengan 0, lintasan berada di titik asal; lalu bergerak ke kanan
menuju (1,0) pada t sama dengan 1, ke atas menuju (1,1) pada t sama dengan 2,
ke kiri menuju (0,1) pada t sama dengan 3, dan akhirnya ke bawah kembali ke
titik asal pada t sama dengan 4.}
\caption{Lintasan $\gamma$ yang menelusuri keliling persegi satuan.\label{fig:squarepath}}
\end{myfigureht}

Lintasan $\gamma$ menelusuri keliling persegi satuan
berlawanan arah jarum jam.  Lihat \figureref{fig:squarepath}.  Lintasan ini
mulus sepotong-sepotong.  Sebagai contoh,
$\gamma|_{[1,2]}(t) = (1,t-1)$ sehingga
$(\gamma|_{[1,2]})'(t) = (0,1) \neq 0$.  Demikian pula untuk ketiga sisi lainnya.
Perhatikan bahwa
$(\gamma|_{[1,2]})'(1) = (0,1)$,
$(\gamma|_{[0,1]})'(1) = (1,0)$, tetapi
$\gamma'(1)$ tidak ada.  Di titik-titik sudut, $\gamma$ 
tidak diferensiabel.
Lintasan $\gamma$ merupakan lintasan tertutup sederhana karena $\gamma|_{[0,4)}$
injektif dan $\gamma(0)=\gamma(4)$.
\end{example}

Definisi lintasan mulus sepotong-sepotong yang kita berikan mengimplikasikan
bahwa lintasan tersebut kontinu (latihan).  Untuk fungsi umum, banyak penulis juga
memperbolehkan berhingga banyak ketakkontinuan ketika menggunakan istilah \emph{mulus
sepotong-sepotong}; karena itu, dapat dikatakan bahwa kita mendefinisikan lintasan mulus
sepotong-sepotong sebagai fungsi yang \emph{\myindex{kontinu dan mulus sepotong-sepotong}}.
Walaupun lintasan mulus saja sudah dapat digunakan, dalam perhitungan lintasan
yang paling mudah dituliskan sering kali mulus sepotong-sepotong.

Pada umumnya, kita lebih tertarik pada citra langsung $\gamma\bigl([a,b]\bigr)$
daripada pada parametrisasi tertentu, walaupun parametrisasi juga
penting sampai taraf tertentu.  Ketika berbicara secara informal tentang lintasan atau kurva,
kita sering memaksudkan himpunan $\gamma\bigl([a,b]\bigr)$, bergantung
pada konteks.


\begin{example}
Syarat $\gamma'(t) \neq 0$ berarti bahwa, secara lokal di sepanjang setiap cabang reguler, citra
$\gamma\bigl([a,b]\bigr)$
tidak memiliki \myquote{sudut} di tempat $\gamma$ mulus.
Tinjau $\gamma \colon \R \to \R^2$ yang didefinisikan oleh
\begin{equation*}
\gamma(t) \coloneqq
\begin{cases}
(t^2,0) & \text{jika } t < 0,\\
(0,t^2) & \text{jika } t \geq 0.
\end{cases}
\end{equation*}
Lihat \figureref{fig:cornersmoothpath}.
Pemeriksaan bahwa $\gamma$ diferensiabel secara kontinu diserahkan kepada pembaca,
namun citra $\gamma(\R) = \bigl\{ (x,y) \in \R^2 : (x,y) =
(s,0) \text{ atau } (x,y) = (0,s) \text{ untuk suatu } s \geq 0 \bigr\}$ memiliki
\myquote{sudut} di titik asal.  Hal ini terjadi karena $\gamma'(0) = (0,0)$.
Terdapat contoh yang lebih rumit, misalnya yang memiliki sudut dalam jumlah tak berhingga;
lihat latihan.
\begin{myfigureht}
\myincludegraphics{cornersmoothpath}{%
Diagram lintasan pada bidang.  Pada t sama dengan minus 1, lintasan berada di
(1,0) pada sumbu x positif, lalu bergerak ke kiri sepanjang sumbu x, melewati
titik bertanda t sama dengan minus satu per dua, hingga mencapai titik asal
pada t sama dengan 0.  Di sana lintasan berbelok ke atas sepanjang sumbu y
positif, melewati t sama dengan satu per dua, lalu mencapai titik bertanda
t sama dengan 1.}
\caption{Lintasan \myquote{mulus} yang memiliki sudut jika turunan yang bernilai nol diizinkan.  Titik-titik
yang bersesuaian dengan beberapa nilai
$t$ ditandai dengan noktah.\label{fig:cornersmoothpath}}
\end{myfigureht}
\end{example}

Pada setiap ruas mulus, syarat $\gamma'(t) \neq 0$ juga di titik-titik ujung
subinterval menjamin bahwa ruas tersebut tidak memiliki sudut dan berakhir secara
reguler; artinya, ruas itu dapat diperpanjang sedikit melampaui titik-titik ujung
sambil tetap reguler jika diperlukan.  Sekali lagi, lihat latihan.

\begin{example}
Untuk $a<b$, grafik fungsi $f \colon [a,b] \to \R$ yang diferensiabel secara kontinu
merupakan lintasan mulus.  Definisikan $\gamma \colon [a,b] \to \R^2$ dengan
\begin{equation*}
\gamma(t) \coloneqq \bigl(t,f(t)\bigr) .
\end{equation*}
Maka $\gamma'(t) = \bigl( 1 , f'(t) \bigr)$, yang tidak pernah nol,
dan $\gamma\bigl([a,b]\bigr)$ merupakan grafik $f$.

Ada cara lain untuk memparametrisasikan lintasan tersebut.  Artinya, terdapat
lintasan-lintasan berbeda dengan citra yang sama.
Fungsi $t \mapsto (1-t)a+tb$ memetakan interval $[0,1]$ ke $[a,b]$.
Definisikan
$\alpha \colon [0,1] \to \R^2$ dengan
\begin{equation*}
\alpha(t) \coloneqq \bigl((1-t)a+tb,f((1-t)a+tb)\bigr) .
\end{equation*}
Maka
$\alpha'(t) = \bigl( b-a ,~ (b-a)f'((1-t)a+tb) \bigr)$, yang tidak pernah nol.
Sebagai himpunan, $\alpha\bigl([0,1]\bigr) = \gamma\bigl([a,b]\bigr)
= \bigl\{ (x,y) \in \R^2 : x \in [a,b] \text{ dan } f(x) = y \bigr\}$,
yang tidak lain adalah grafik $f$.
\end{example}

Contoh terakhir mengantar kita pada suatu definisi.

\begin{defn}
Misalkan $\gamma \colon [a,b] \to \R^n$ lintasan mulus dan
$h \colon [c,d] \to [a,b]$ fungsi bijektif yang diferensiabel secara kontinu
sedemikian sehingga $h'(t) \neq 0$ untuk semua $t \in [c,d]$.  Maka
komposisi
$\gamma \circ h$ disebut
\emph{\myindex{reparametrisasi mulus}}\index{reparametrisasi}
dari $\gamma$.

Misalkan $\gamma \colon [a,b] \to \R^n$ lintasan mulus sepotong-sepotong dan
$h \colon [c,d] \to [a,b]$ fungsi bijektif yang mulus sepotong-sepotong,
dengan semua limit sepihak dari $h'$ pada titik-titik partisi ada dan tak nol.
Komposisi
$\gamma \circ h$ disebut
\emph{\myindex{reparametrisasi mulus sepotong-sepotong}} dari $\gamma$.

Jika $h$ naik secara ketat, dikatakan bahwa $h$ 
\emph{\myindex{mempertahankan orientasi}}.  Jika $h$ turun secara ketat, dikatakan bahwa
$h$ \emph{\myindex{membalik orientasi}}.
\end{defn}

Reparametrisasi merupakan lintasan lain dengan citra yang sama.  Artinya,
$(\gamma \circ h)\bigl([c,d]\bigr) =
\gamma \bigl([a,b]\bigr)$.

Syarat pada $h$ yang mulus sepotong-sepotong berarti bahwa
terdapat suatu partisi $t_0 = c < t_1 < t_2 < \cdots < t_k = d$
sedemikian sehingga $h|_{[t_{j-1},t_j]}$ diferensiabel secara kontinu
dan $\bigl(h|_{[t_{j-1},t_j]}\bigr)'(t) \neq 0$
untuk semua $t \in [t_{j-1},t_j]$.
Karena $h$ bijektif, fungsi tersebut naik ketat atau
turun ketat.  Jadi, $\bigl(h|_{[t_{j-1},t_j]}\bigr)'(t) > 0$
untuk semua $t$, atau $\bigl(h|_{[t_{j-1},t_j]}\bigr)'(t) < 0$ untuk semua $t$.

\begin{prop} \label{prop:reparamapiecewisesmooth}
Jika $\gamma \colon [a,b] \to \R^n$ merupakan lintasan mulus sepotong-sepotong
dan $\gamma \circ h \colon [c,d] \to \R^n$ merupakan
reparametrisasi mulus sepotong-sepotong, maka $\gamma \circ h$
merupakan lintasan mulus sepotong-sepotong.
\end{prop}

\begin{proof}
Andaikan $h$ mempertahankan orientasi; artinya, $h$ naik ketat.
Jika $h \colon [c,d] \to [a,b]$ memberikan reparametrisasi mulus sepotong-sepotong,
maka untuk suatu partisi
$r_0 = c < r_1 < r_2 < \cdots < r_\ell = d$, pembatasan
$h|_{[r_{j-1},r_j]}$ diferensiabel secara kontinu dengan turunan
positif.

Kita perhalus partisi dari definisi kemulusan sepotong-sepotong untuk $\gamma$
dengan menambahkan titik-titik
$\{ h(r_0), h(r_1), h(r_2), \ldots, h(r_\ell) \}$, lalu urutkan dan namai ulang
partisi itu sebagai $t_0 = a < t_1 < t_2 < \cdots < t_k = b$.
Misalkan $s_j \coloneqq h^{-1}(t_j)$.  Maka
$s_0 = c < s_1 < s_2 < \cdots < s_k = d$
merupakan partisi yang memuat semua titik $r_i$, sehingga merupakan penghalusan
dari partisi $r_0 < r_1 < \cdots < r_\ell$.
Jika $\tau \in [s_{j-1},s_j]$, maka $h(\tau) \in [t_{j-1},t_j]$
karena $h(s_{j-1}) = t_{j-1}$,
$h(s_{j}) = t_j$, dan
$h$ naik ketat.
Selain itu, $h|_{[s_{j-1},s_j]}$ diferensiabel secara kontinu, dan
$\gamma|_{[t_{j-1},t_j]}$ juga diferensiabel secara kontinu.
Dengan demikian,
\begin{equation*}
(\gamma \circ h)|_{[s_{j-1},s_{j}]} (\tau)
=
\gamma|_{[t_{j-1},t_{j}]} \bigl( h|_{[s_{j-1},s_j]}(\tau) \bigr) .
\end{equation*}
Oleh karena itu, fungsi 
$(\gamma \circ h)|_{[s_{j-1},s_{j}]}$
diferensiabel secara kontinu dan, menurut aturan rantai,
\begin{equation*}
\bigl( (\gamma \circ h)|_{[s_{j-1},s_{j}]} \bigr) ' (\tau)
=
\bigl( \gamma|_{[t_{j-1},t_{j}]} \bigr)' \bigl( h(\tau) \bigr)
\bigl(h|_{[s_{j-1},s_j]}\bigr)'(\tau) \neq 0 .
\end{equation*}
Akibatnya, $\gamma \circ h$ merupakan lintasan mulus sepotong-sepotong.
Pembuktian untuk $h$ yang membalik orientasi diserahkan sebagai latihan.
\end{proof}

Jika dua lintasan sederhana yang mulus sepotong-sepotong dan reguler serta tidak tertutup memiliki citra yang sama,
pembuktian bahwa terdapat reparametrisasi diserahkan sebagai latihan.  Pernyataan
yang sama berlaku untuk lintasan tertutup sederhana yang mulus sepotong-sepotong
dan reguler apabila titik pangkal yang dipilih juga sama; jika titik pangkalnya berbeda, diperlukan reparametrisasi siklik
pada lingkaran parameter.  Di sinilah asumsi bahwa turunan pada setiap ruas mulus
dari kedua lintasan tidak pernah nol menjadi penting.

\subsection{Integral lintasan bentuk-1}

\begin{defn}
Gunakan $(x_1,x_2,\ldots,x_n) \in \R^n$ sebagai koordinat kita.
Diberikan $n$ fungsi kontinu bernilai riil
$\omega_1,\omega_2,\ldots,\omega_n$ yang didefinisikan pada himpunan $S \subset \R^n$,
kita mendefinisikan \emph{\myindex{bentuk-1}}\index{bentuk diferensial derajat satu}
sebagai objek berbentuk
\glsadd{not:oneform}
\begin{equation*}
\omega = \omega_1 \,dx_1 + \omega_2 \,dx_2 + \cdots + \omega_n \,dx_n .
\end{equation*}
Kita dapat merepresentasikan $\omega$ sebagai fungsi kontinu dari $S$ ke $\R^n$,
tetapi lebih baik memandangnya sebagai jenis objek yang berbeda.
\end{defn}

\begin{example}
\begin{equation*}
\omega(x,y) \coloneqq \frac{-y}{x^2+y^2} \,dx + \frac{x}{x^2+y^2} \,dy
\end{equation*}
merupakan bentuk-1 yang didefinisikan pada $\R^2 \setminus \{ (0,0) \}$.
\end{example}

\begin{defn}
Misalkan $\gamma \colon [a,b] \to \R^n$ lintasan mulus
dan misalkan
\begin{equation*}
\omega = \omega_1 \,dx_1 + \omega_2 \,dx_2 + \cdots + \omega_n \,dx_n
\end{equation*}
merupakan bentuk-1 yang didefinisikan pada citra langsung $\gamma\bigl([a,b]\bigr)$.
Tuliskan $\gamma = (\gamma_1,\gamma_2,\ldots,\gamma_n)$.
Definisikan:
\begin{equation*}
\begin{split}
\int_{\gamma} \omega
& \coloneqq
\int_a^b 
\Bigl(
\omega_1\bigl(\gamma(t)\bigr) \gamma_1'(t) +
\omega_2\bigl(\gamma(t)\bigr) \gamma_2'(t) + \cdots +
\omega_n\bigl(\gamma(t)\bigr) \gamma_n'(t) \Bigr) dt
\\
&\phantom{:}=
\int_a^b 
\left(
\sum_{j=1}^n
\omega_j\bigl(\gamma(t)\bigr) \gamma_j'(t) \right) dt .
\end{split}
\end{equation*}
Untuk mengingat definisi ini, perhatikan bahwa $x_j$ adalah $\gamma_j(t)$; sebagai
kaidah substitusi, $dx_j$ menjadi $\gamma_j'(t) \, dt$.

Jika $\gamma$ mulus sepotong-sepotong, ambil partisi yang bersesuaian
$t_0 = a < t_1 < t_2 < \ldots < t_k = b$, dan andaikan partisi tersebut
minimal
%\footnote{Pembatasan ini sebenarnya tidak mutlak diperlukan, tetapi 
%membuat integral langsung terdefinisi dengan baik.}
dalam arti bahwa $\gamma$ tidak diferensiabel
di setiap $t_1,t_2,\ldots,t_{k-1}$.  Karena setiap $\gamma|_{[t_{j-1},t_j]}$
merupakan lintasan mulus, definisikan
\begin{equation*}
\int_{\gamma} \omega
\coloneqq
\int_{\gamma|_{[t_0,t_1]}} \omega
\,
+
\,
\int_{\gamma|_{[t_1,t_2]}} \omega
\,
+ \, \cdots \, + \,
\int_{\gamma|_{[t_{k-1},t_k]}} \omega .
\end{equation*}
\end{defn}

Notasi integral lintasan bentuk-1 dapat dipahami melalui rumus
\myquote{substitusi $u$} yang dikenal dari kalkulus.
Secara agak informal, sebagai kaidah substitusi,
jika $x_j(t) = \gamma_j(t)$, maka $dx_j$ menjadi $\gamma_j'(t) \, dt$.

Lintasan dapat dipotong menjadi sublintasan atau disambungkan.  Pembuktiannya merupakan penerapan langsung
dari aditivitas integral Riemann dan diserahkan sebagai latihan.
Proposisi berikut membenarkan cara kita mendefinisikan integral pada lintasan
mulus sepotong-sepotong.  Proposisi itu juga menunjukkan bahwa kita dapat
menggunakan sembarang partisi yang sesuai, bukan hanya partisi minimal dalam definisi.

\begin{prop} \label{mv:prop:pathconcat}
Misalkan $\gamma \colon [a,c] \to \R^n$ lintasan mulus sepotong-sepotong
dan $b \in (a,c)$.
Definisikan lintasan-lintasan mulus sepotong-sepotong
$\alpha \coloneqq \gamma|_{[a,b]}$ dan
$\beta \coloneqq \gamma|_{[b,c]}$.
Misalkan $\omega$ bentuk-1 yang didefinisikan pada
$\gamma\bigl([a,c]\bigr)$.  Maka
\begin{equation*}
\int_{\gamma} \omega =
\int_{\alpha} \omega +
\int_{\beta} \omega .
\end{equation*}
\end{prop}


\begin{example} \label{example:mv:irrotoneformint}
Misalkan bentuk-1 $\omega$ dan lintasan $\gamma \colon [0,2\pi] \to \R^2$ didefinisikan oleh
\begin{equation*}
\omega(x,y) \coloneqq \frac{-y}{x^2+y^2} \,dx + \frac{x}{x^2+y^2} \,dy,
\qquad
\gamma(t) \coloneqq \bigl(\cos(t),\sin(t)\bigr) .
\end{equation*}
Maka
\begin{equation*}
\begin{split}
\int_{\gamma} \omega
& =
\int_0^{2\pi}
\Biggl(
\frac{-\sin(t)}{{\bigl(\cos(t)\bigr)}^2+{\bigl(\sin(t)\bigr)}^2}
\bigl(-\sin(t)\bigr)
%mbxlatex \\
%mbxlatex & \phantom{=\int_0^{2\pi}\Biggl(~}
+
\frac{\cos(t)}{{\bigl(\cos(t)\bigr)}^2+{\bigl(\sin(t)\bigr)}^2}
\bigl(\cos(t)\bigr)
\Biggr) \, dt
\\
& =
\int_0^{2\pi}
1 \, dt
= 2\pi .
\end{split}
\end{equation*}
Selanjutnya, parametrisasikan ulang kurva yang sama menggunakan
$\alpha \colon [0,1] \to \R^2$ yang didefinisikan oleh $\alpha(t) \coloneqq \bigl(\cos(2\pi
t),\sin(2 \pi t)\bigr)$.  Dengan $h(t)=2\pi t$, kita mempunyai
$\alpha=\gamma\circ h$, sehingga $\alpha$ merupakan reparametrisasi mulus dari $\gamma$.
Maka
\begin{equation*}
\begin{split}
\int_{\alpha} \omega
& =
\int_0^{1}
\Biggl(
\frac{-\sin(2\pi t)}{{\bigl(\cos(2\pi t)\bigr)}^2+{\bigl(\sin(2\pi t)\bigr)}^2}
\bigl(-2\pi \sin(2\pi t)\bigr)
\\
& \phantom{=\int_0^1\Biggl(~}
+
\frac{\cos(2 \pi t)}{{\bigl(\cos(2 \pi t)\bigr)}^2+{\bigl(\sin(2 \pi t)\bigr)}^2}
\bigl(2 \pi \cos(2 \pi t)\bigr)
\Biggr) \, dt
\\
& =
\int_0^{1}
2\pi \, dt
= 2\pi .
\end{split}
\end{equation*}
Terakhir, lakukan reparametrisasi dengan $\beta \colon [0,2\pi] \to \R^2$
yang didefinisikan oleh $\beta(t) \coloneqq \bigl(\cos(-t),\sin(-t)\bigr)$.
Secara eksplisit, $\beta=\gamma\circ h$ untuk $h(t)=2\pi-t$.  Maka
\begin{equation*}
\begin{split}
\int_{\beta} \omega
& =
\int_0^{2\pi}
\Biggl(
\frac{-\sin(-t)}{{\bigl(\cos(-t)\bigr)}^2+{\bigl(\sin(-t)\bigr)}^2}
\bigl(\sin(-t)\bigr)
%mbxlatex \\
%mbxlatex & \phantom{=\int_0^{2\pi}\Biggl(~}
+
\frac{\cos(-t)}{{\bigl(\cos(-t)\bigr)}^2+{\bigl(\sin(-t)\bigr)}^2}
\bigl(-\cos(-t)\bigr)
\Biggr) \, dt
\\
& =
\int_0^{2\pi}
(-1) \, dt
= -2\pi .
\end{split}
\end{equation*}
Lintasan $\alpha$ merupakan reparametrisasi dari $\gamma$ yang mempertahankan
orientasi, dan kedua integralnya sama.  Lintasan $\beta$
merupakan reparametrisasi dari $\gamma$ yang membalik orientasi, dan integralnya
merupakan negatif dari integral semula.  Lihat \figureref{fig:circlepathrepar}.
\begin{myfigureht}
\myincludegraphics{circlepathrepar}{%
Lintasan melingkari titik asal.  Panah menunjukkan arah berlawanan jarum jam
untuk gamma dan alfa, sedangkan beta bergerak searah jarum jam.
Titik pada sumbu x positif merupakan titik awal bagi $\gamma$ dan $\alpha$ serta
titik akhir bagi $\beta$, dan diberi label gamma dari 0 sama dengan alfa dari 0
sama dengan beta dari 2 pi.  Titik yang sama juga memenuhi gamma dari 2 pi sama
dengan alfa dari 1 sama dengan beta dari 0.
Titik berikutnya di kuadran pertama, di pertengahan menuju sumbu y positif,
diberi label gamma dari pi per 4 sama dengan alfa dari satu per delapan sama dengan
beta dari 7 pi per 4.  Titik pada sumbu y positif diberi label gamma dari pi per 2
sama dengan alfa dari satu per empat sama dengan beta dari 3 pi per 2.  Titik pada
sumbu x negatif diberi label gamma dari pi sama dengan alfa dari satu per dua sama
dengan beta dari pi.  Titik pada sumbu y negatif diberi label gamma dari 3 pi per 2
sama dengan alfa dari tiga per empat sama dengan beta dari pi per 2.}
\caption{Lintasan lingkaran yang direparametrisasi dengan dua cara berbeda. Panah
menunjukkan orientasi $\gamma$ dan $\alpha$. Lintasan $\beta$ menelusuri
lingkaran dalam arah yang berlawanan.\label{fig:circlepathrepar}}
\end{myfigureht}
\end{example}

Contoh sebelumnya bukan suatu kebetulan.
Nilai integral lintasan tidak bergantung pada pilihan reparametrisasi bijektif,
kecuali melalui orientasi: pembalikan orientasi mengubah tandanya.  Parametrisasi
sebarang dengan citra yang sama tidak termasuk dalam pernyataan ini karena banyaknya penelusuran juga berpengaruh.

\begin{prop} \label{mv:prop:pathintrepararam}
Misalkan $\gamma \colon [a,b] \to \R^n$ lintasan mulus sepotong-sepotong dan
$\gamma \circ h \colon [c,d] \to \R^n$ reparametrisasi mulus sepotong-sepotong.
Andaikan $\omega$ bentuk-1 yang didefinisikan pada himpunan $\gamma\bigl([a,b]\bigr)$.  Maka
\begin{equation*}
\int_{\gamma \circ h} \omega =
\begin{cases}
\int_{\gamma} \omega  & \text{jika } h \text{ mempertahankan orientasi,}\\
-\int_{\gamma} \omega & \text{jika } h \text{ membalik orientasi.}
\end{cases}
\end{equation*}
\end{prop}

\begin{proof}
Mula-mula, andaikan $\gamma$ dan $h$ keduanya mulus.
Tuliskan $\omega = \omega_1 \, dx_1 + \omega_2 \, dx_2 + \cdots +
\omega_n \, dx_n$.
Andaikan $h$ mempertahankan orientasi.  Gunakan
rumus penggantian variabel untuk integral Riemann:
\begin{equation*}
\begin{split}
\int_{\gamma} \omega
& =
\int_a^b 
\left(
\sum_{j=1}^n
\omega_j\bigl(\gamma(t)\bigr) \gamma_j'(t)
\right) dt
\\
& =
\int_c^d 
\left(
\sum_{j=1}^n
\omega_j\Bigl(\gamma\bigl(h(\tau)\bigr)\Bigr) \gamma_j'\bigl(h(\tau)\bigr)
\right) h'(\tau) \, d\tau
\\
& =
\int_c^d 
\left(
\sum_{j=1}^n
\omega_j\Bigl(\gamma\bigl(h(\tau)\bigr)\Bigr) (\gamma_j \circ h)'(\tau)
\right) d\tau
=
\int_{\gamma \circ h} \omega .
\end{split}
\end{equation*}
Jika $h$ membalik orientasi, penggantian variabel menukar urutan batas integrasi
dan menghasilkan tanda minus.
Rincian ini, beserta penyelesaian pembuktian untuk lintasan mulus
sepotong-sepotong, diserahkan kepada pembaca dalam \exerciseref{mv:exercise:pathpiece}.
\end{proof}

Berdasarkan proposisi ini (dan latihan), jika $\Gamma
\subset \R^n$ merupakan citra lintasan mulus sepotong-sepotong yang sederhana,
yaitu $\gamma\bigl([a,b]\bigr)$, maka asalkan orientasinya ditentukan, yaitu arah
penelusuran kurva, kita dapat menuliskan
\begin{equation*}
\int_{\Gamma} \omega ,
\end{equation*}
tanpa menyebutkan $\gamma$ tertentu.
Selain itu, untuk lintasan tertutup sederhana, titik awal parametrisasi juga
tidak memengaruhi integral: lintasan dapat dipotong di titik awal baru lalu
disambungkan kembali.  Lihat latihan.

Ingat bahwa \emph{sederhana} berarti $\gamma$ injektif, kecuali bahwa
$\gamma(a)=\gamma(b)$ diperbolehkan untuk lintasan tertutup sederhana; khususnya,
$\gamma$ injektif ketika dibatasi pada $[a,b)$.
Kita dapat sedikit memperlemah syarat bahwa lintasan harus sederhana.
Sebagai contoh, cukup andaikan terdapat himpunan berhingga $E \subset [a,b]$
sedemikian sehingga
$\gamma|_{[a,b]\setminus E}$ injektif.
%Artinya, terdapat himpunan berhingga $E \subset [a,b]$
%dan $\gamma|_{[a,b]\setminus E}$ injektif.
Lihat \exerciseref{mv:exercise:curveintegral}.  Namun, syarat tersebut tidak dapat
dihilangkan sepenuhnya, seperti diperlihatkan oleh contoh berikut.

\begin{example}
Ambil $\gamma \colon [0,2\pi] \to \R^2$ yang diberikan oleh $\gamma(t) \coloneqq
\bigl(\cos(t),\sin(t)\bigr)$, dan
$\beta \colon [0,2\pi] \to \R^2$ dengan $\beta(t) \coloneqq
\bigl(\cos(2t),\sin(2t)\bigr)$.  Perhatikan bahwa
$\gamma\bigl([0,2\pi]\bigr) = \beta\bigl([0,2\pi]\bigr)$; kedua lintasan menelusuri
kurva yang sama, yaitu lingkaran satuan.  Akan tetapi, $\gamma$ mengelilingi lingkaran
satuan satu kali berlawanan arah jarum jam, sedangkan $\beta$ mengelilingi
lingkaran satuan dua kali (dalam arah yang sama). 
Lihat \figureref{fig:circlepathrepar2}.
\begin{myfigureht}
\myincludegraphics{circlepathrepar2}{%
Lintasan mengelilingi titik asal berlawanan arah jarum jam.  Gamma menelusuri
lingkaran satu kali, sedangkan beta menelusurinya dua kali.  Titik awal pada sumbu x
positif diberi label gamma dari 0 sama dengan beta dari 0 sama dengan beta dari pi.
Titik di kuadran pertama, di pertengahan menuju sumbu y positif, diberi label
gamma dari pi per 4 sama dengan beta dari pi per 8 sama dengan beta dari 9 pi per 8.
Titik pada sumbu y positif diberi label gamma dari pi per 2 sama dengan beta dari
pi per 4 sama dengan beta dari 5 pi per 4.  Titik pada sumbu x negatif diberi label
gamma dari pi sama dengan beta dari pi per 2 sama dengan beta dari 3 pi per 2.
Titik pada sumbu y negatif diberi label gamma dari 3 pi per 2 sama dengan beta dari
3 pi per 4 sama dengan beta dari 7 pi per 4.  Titik akhir yang sama dengan titik awal
diberi label gamma dari 2 pi sama dengan beta dari 2 pi.}
\caption{Lintasan lingkaran yang ditelusuri satu kali oleh
$\gamma \colon [0,2\pi] \to \R^2$
dan dua kali oleh
$\beta \colon [0,2\pi] \to \R^2$.\label{fig:circlepathrepar2}}
\end{myfigureht}

Hitung
\begin{align*}
& \int_{\gamma} -y\, dx + x\,dy
=
\int_0^{2\pi}
\Bigl( \bigl(-\sin(t) \bigr) \bigl(-\sin(t) \bigr) + \cos(t) \cos(t) \Bigr) dt
=
2 \pi,\\
& \int_{\beta} -y\, dx + x\,dy
=
\int_0^{2\pi}
\Bigl( \bigl(-\sin(2t) \bigr) \bigl(-2\sin(2t) \bigr) + \cos(2t)
\bigl(2\cos(2t)\bigr) \Bigr) dt
=
4 \pi.
\end{align*}
\end{example}

Kadang-kadang berguna untuk mendefinisikan integral lintasan bagi suatu pemetaan
$\gamma \colon [a,b] \to \R^n$ yang diferensiabel secara kontinu, tetapi tidak harus
merupakan lintasan reguler menurut definisi kita.  Andaikan $\omega$ didefinisikan
pada $\gamma\bigl([a,b]\bigr)$.  Definisikan
\glsadd{not:pathintegralomega}
\begin{equation*}
\int_{\gamma} \omega \coloneqq \int_a^b
\left(
\sum_{j=1}^n
\omega_j\bigl(\gamma(t)\bigr) \gamma_j'(t)
\right) dt 
\end{equation*}
untuk setiap $\gamma$ yang diferensiabel secara kontinu.  Kasus yang
muncul secara alami adalah ketika $\gamma$ konstan.  Dalam hal ini,
$\gamma'(t) = 0$ untuk semua $t$, dan $\gamma\bigl([a,b]\bigr)$ merupakan satu
titik, yang kita pandang sebagai \myquote{kurva} dengan panjang nol.  Maka
$\int_{\gamma} \omega = 0$ untuk setiap $\omega$.

\subsection{Integral lintasan suatu fungsi}

Selanjutnya, kita mengintegralkan suatu fungsi terhadap
\emph{\myindex{ukuran panjang busur}} $ds$.  Untuk fungsi tak negatif, gambaran geometrisnya
adalah luas di bawah grafik fungsi di atas suatu lintasan.
Bayangkan sebuah pagar tegak di atas $\gamma$, dengan tinggi yang diberikan oleh fungsi.
Integralnya adalah luas pagar tersebut.  Untuk fungsi bernilai riil umum, integral ini
bertanda; luas geometris pagar diperoleh dengan mengintegralkan nilai mutlak fungsi.
Lihat \figureref{fig:fenceintegral}.
\begin{myfigureht}
\myincludegraphics{fenceintegral}{%
Diagram dalam ruang xyz.  Kurva gamma digambar sebagai garis tebal pada
bidang xy.  Sebuah permukaan vertikal berarsir seperti pagar berdiri di atas
kurva tersebut, dengan tinggi yang berbeda-beda di atas titik-titik yang berbeda.
Untuk fungsi tak negatif, luas permukaan berarsir menggambarkan integral.}
\caption{Lintasan $\gamma \colon [a,b] \to \R^2$ pada bidang $xy$ (kurva tebal), dan fungsi
$z=f(x,y)$ yang digambarkan di atasnya dalam arah $z$.  Untuk $f\geq0$, integral merupakan
luas berarsir yang ditampilkan.\label{fig:fenceintegral}}
\end{myfigureht}

\begin{defn}
Andaikan $\gamma \colon [a,b] \to \R^n$ lintasan mulus, dan
$f \colon \gamma\bigl([a,b]\bigr) \to \R$ fungsi kontinu.  Definisikan
\glsadd{not:lineintegralf}
\begin{equation*}
\int_{\gamma} f \,ds \coloneqq
\int_a^b f\bigl( \gamma(t) \bigr) \bnorm{\gamma'(t)} \, dt .
\end{equation*}
Untuk menegaskan variabelnya, kita dapat menggunakan
\begin{equation*}
\int_{\gamma} f(x) \,ds(x) \coloneqq \int_{\gamma} f \,ds .
\end{equation*}

Definisi untuk lintasan mulus sepotong-sepotong serupa dengan definisi sebelumnya
dan diserahkan kepada pembaca.
\end{defn}

Integral lintasan suatu fungsi juga tidak bergantung pada pilihan reparametrisasi
bijektif, dan dalam hal ini orientasi tidak berpengaruh.

\begin{prop} \label{mv:prop:lineintrepararam}
Misalkan $\gamma \colon [a,b] \to \R^n$ lintasan mulus sepotong-sepotong dan
$\gamma \circ h \colon [c,d] \to \R^n$ reparametrisasi mulus sepotong-sepotong.
Andaikan $f$ fungsi kontinu yang didefinisikan pada himpunan
$\gamma\bigl([a,b]\bigr)$.  Maka
\begin{equation*}
\int_{\gamma \circ h} f\, ds = \int_{\gamma} f\, ds .
\end{equation*}
\end{prop}

\begin{proof}
Andaikan $h$ mempertahankan orientasi serta $\gamma$ dan $h$
keduanya mulus.  Maka 
\begin{equation*}
\begin{split}
\int_{\gamma} f \, ds
& =
\int_a^b 
f\bigl(\gamma(t)\bigr) \bnorm{\gamma'(t)} \, dt
\\
& =
\int_c^d 
f\Bigl(\gamma\bigl(h(\tau)\bigr)\Bigr)
\bnorm{\gamma'\bigl(h(\tau)\bigr)} h'(\tau) \, d\tau
\\
& =
\int_c^d 
f\Bigl(\gamma\bigl(h(\tau)\bigr)\Bigr)
\bnorm{\gamma'\bigl(h(\tau)\bigr) h'(\tau)} \, d\tau
\\
& =
\int_c^d 
f\bigl((\gamma \circ h)(\tau)\bigr) \bnorm{(\gamma \circ h)'(\tau)} \, d\tau
\\
& = 
\int_{\gamma \circ h} f \, ds .
\end{split}
\end{equation*}
Jika $h$ membalik orientasi, penggantian variabel menukar urutan batas
integrasi.  Karena $h'(\tau)<0$, berlaku
$\bnorm{\gamma'(h(\tau))h'(\tau)}=\bnorm{\gamma'(h(\tau))}\sabs{h'(\tau)}$;
tanda minus dari pembalikan batas tepat mengimbangi
$\sabs{h'(\tau)}=-h'(\tau)$.
Rincian ini, beserta penyelesaian pembuktian untuk lintasan mulus
sepotong-sepotong, diserahkan kepada pembaca dalam \exerciseref{mv:exercise:linepiece}.
\end{proof}

Seperti sebelumnya, berdasarkan proposisi ini (dan latihan),
jika $\gamma$ sederhana, pilihan reparametrisasi bijektif tidak berpengaruh.
Untuk lintasan tertutup sederhana, perubahan titik awal ditangani dengan memotong
dan menyambungkan lintasan secara siklik, lalu menggunakan aditivitas.
Oleh karena itu, jika $\Gamma = \gamma\bigl( [a,b] \bigr)$, kita cukup
menuliskan
\begin{equation*}
\int_\Gamma f\, ds .
\end{equation*}
Dalam hal ini kita juga tidak perlu memperhatikan orientasi; kedua arah
menghasilkan integral yang sama.

\begin{example}
Misalkan $f(x,y) \coloneqq x$.  Misalkan $C \subset \R^2$ bagian lingkaran satuan
pada setengah bidang $x \geq 0$.  Kita hendak menghitung
\begin{equation*}
\int_C f \, ds .
\end{equation*}
Parametrisasikan kurva $C$ dengan $\gamma \colon
[\nicefrac{-\pi}{2},\nicefrac{\pi}{2}] \to \R^2$ yang didefinisikan oleh
$\gamma(t) \coloneqq \bigl(\cos(t),\sin(t)\bigr)$.
Maka $\gamma'(t) = \bigl(-\sin(t),\cos(t)\bigr)$, dan
\begin{equation*}
\int_C f \, ds =
\int_\gamma f \, ds
=
\int_{-\pi/2}^{\pi/2} \cos(t) \sqrt{ {\bigl(-\sin(t)\bigr)}^2 +  
{\bigl(\cos(t)\bigr)}^2 } \, dt
=
\int_{-\pi/2}^{\pi/2} \cos(t) \, dt = 2.
\end{equation*}
\end{example}

\begin{defn}
Andaikan $\Gamma \subset \R^n$ diparametrisasikan oleh lintasan sederhana
yang mulus sepotong-sepotong, $\gamma \colon [a,b] \to \R^n$; artinya,
$\gamma\bigl( [a,b] \bigr) = \Gamma$.  Kita mendefinisikan
\emph{\myindex{panjang}}\index{panjang kurva} dengan
\begin{equation*}
\ell(\Gamma) \coloneqq \int_{\Gamma} ds = \int_{\gamma} ds .
\end{equation*}
\end{defn}

Jika $\gamma$ mulus,
\begin{equation*}
\ell(\Gamma) = 
\int_a^b
\bnorm{\gamma'(t)}\, dt .
\end{equation*}
Perlu disebutkan bahwa lazim untuk menuliskan
$\int_a^b
\bnorm{\gamma'(t)}\, dt$ meskipun lintasannya hanya mulus sepotong-sepotong.
$\gamma'(t)$ ada dan kontinu kecuali mungkin pada berhingga banyak titik partisi.
Pada titik-titik tersebut, perluas fungsi yang dinotasikan dengan $\gamma'$ dengan
menetapkan nilai sembarang, misalnya nol.  Fungsi $\bnorm{\gamma'(t)}$ yang
dihasilkan terbatas dan kontinu kecuali di berhingga banyak titik, sehingga integral
Riemann tersebut ada.

\begin{example}
Misalkan $x,y \in \R^n$ dua titik berbeda dan tuliskan $[x,y]$ untuk ruas garis lurus
di antara $x$ dan $y$.  Parametrisasikan
$[x,y]$ dengan $\gamma(t) \coloneqq (1-t)x + ty$ untuk $t$ dari $0$ hingga $1$.
Lihat \figureref{fig:straightpath}.
Maka $\gamma'(t) = y-x$, sehingga
\begin{equation*}
\ell\bigl([x,y]\bigr)
=
\int_{[x,y]} ds
=
\int_0^1 \snorm{y-x} \, dt
=
\snorm{y-x} .
\end{equation*}
Panjang $[x,y]$ adalah jarak Euklides standar antara $x$ dan $y$,
yang membenarkan penamaan tersebut.  Jika $x=y$, rumus yang sama tetap berlaku
melalui konvensi kurva konstan dengan panjang nol.
\begin{myfigureht}
\myincludegraphics{straightpath}{%
Diagram lintasan garis lurus dari titik x menuju titik y, dengan arah
ditunjukkan oleh sebuah panah.  Lintasan diberi label [x,y].  Titik x juga
diberi label t sama dengan 0, dan titik y diberi label t sama dengan 1.}
\caption{Lintasan lurus antara $x$ dan $y$ yang diparametrisasikan
oleh $(1-t)x + ty$.\label{fig:straightpath}}
\end{myfigureht}
\end{example}

Untuk $r>0$, lintasan sederhana dan mulus sepotong-sepotong
$\gamma \colon [0,r] \to \R^n$ disebut
\emph{\myindex{parametrisasi panjang busur}} jika
untuk semua $t \in [0,r]$ berlaku
\begin{equation*}
\ell\bigl( \gamma\bigl([0,t]\bigr) \bigr) = t .
\end{equation*}
Jika $\gamma$ mulus, maka
\begin{equation*}
\int_0^t d\tau =
t = \ell\bigl( \gamma\bigl([0,t]\bigr) \bigr) =
\int_0^t
\bnorm{\gamma'(\tau)}
\, d\tau
\end{equation*}
untuk semua $t$.
Teorema dasar kalkulus kemudian menyatakan
bahwa $\bnorm{\gamma'(t)} = 1$ untuk semua $t$.
Demikian pula, untuk $\gamma$ yang mulus sepotong-sepotong, diperoleh
$\bnorm{\gamma'(t)} = 1$ untuk semua $t$ tempat turunannya ada.
Pada titik-titik partisi, turunan sepihak pada setiap ruas yang bersebelahan
juga mempunyai norma satu.
Bayangkan parametrisasi semacam ini sebagai gerak di sepanjang kurva dengan kelajuan
1.  Jika $\gamma \colon [0,r] \to \R^n$ merupakan parametrisasi panjang busur, lazim
menggunakan $s$ sebagai variabel, sehingga $\int_\gamma f \,ds
= \int_0^r f\bigl(\gamma(s)\bigr) \,ds$.

\subsection{Latihan}

\begin{exercise}
Buktikan bahwa jika $\varphi \colon [a,b] \to \R^n$ merupakan lintasan mulus
sepotong-sepotong menurut definisi kita, maka $\varphi$ merupakan fungsi kontinu.
\end{exercise}

\begin{exercise}
Selesaikan pembuktian \propref{prop:reparamapiecewisesmooth} untuk reparametrisasi
yang membalik orientasi.
\end{exercise}

\begin{exercise}
Buktikan \propref{mv:prop:pathconcat}.
\end{exercise}

\begin{exercise} \label{mv:exercise:pathpiece}
Selesaikan pembuktian \propref{mv:prop:pathintrepararam}
untuk
\begin{enumerate}[a)]
\item
reparametrisasi yang membalik orientasi, dan
\item
lintasan serta reparametrisasi yang mulus
sepotong-sepotong.
\end{enumerate}
\end{exercise}

\begin{exercise} \label{mv:exercise:linepiece}
Selesaikan pembuktian \propref{mv:prop:lineintrepararam}
untuk
\begin{enumerate}[a)]
\item
reparametrisasi yang membalik orientasi, dan
\item
lintasan serta reparametrisasi yang mulus
sepotong-sepotong.
\end{enumerate}
\end{exercise}

\begin{exercise}
Andaikan $\gamma \colon [a,b] \to \R^n$ lintasan mulus sepotong-sepotong dan $f$
fungsi kontinu yang didefinisikan pada citra $\gamma\bigl([a,b]\bigr)$.
Berikan definisi untuk $\int_{\gamma} f \,ds$.
\end{exercise}

\begin{exercise}
Dengan menggunakan definisi secara langsung, hitung:
\begin{enumerate}[a)]
\item
Panjang busur keliling persegi satuan dari
\exampleref{mv:example:unitsquarepath} dengan parametrisasi yang diberikan.
\item
Panjang busur lingkaran satuan dengan parametrisasi
$\gamma \colon [0,1] \to \R^2$, $\gamma(t) \coloneqq \bigl(\cos(2\pi t),\sin(2\pi t)\bigr)$.
\item
Panjang busur lingkaran satuan dengan parametrisasi
$\beta \colon [0,2\pi] \to \R^2$, $\beta(t) \coloneqq \bigl(\cos(t),\sin(t)\bigr)$.
\end{enumerate}
Catatan: Anda boleh menggunakan pengetahuan tentang sinus dan kosinus dari kalkulus.
\end{exercise}

\begin{exercise}
Andaikan $\gamma \colon [0,1] \to \R^n$ lintasan mulus dan
$\omega$ bentuk-1 yang didefinisikan pada citra $\gamma\bigl([0,1]\bigr)$.
Untuk $r \in (0,1]$, definisikan $\gamma_r \colon [0,r] \to \R^n$
sebagai pembatasan $\gamma$ pada $[0,r]$, lalu definisikan
$h(r) \coloneqq \int_{\gamma_r} \omega$ dan tetapkan $h(0)\coloneqq 0$.
Buktikan bahwa $h$ diferensiabel secara kontinu pada $[0,1]$.
\end{exercise}

\begin{exercise}
Andaikan $\gamma \colon [a,b] \to \R^n$ lintasan mulus.
Buktikan bahwa terdapat $\epsilon > 0$ dan fungsi yang diferensiabel secara kontinu
$\widetilde{\gamma} \colon (a-\epsilon,b+\epsilon) \to \R^n$
dengan $\widetilde{\gamma}(t) = \gamma(t)$ untuk semua $t \in [a,b]$
dan $\widetilde{\gamma}'(t) \neq 0$ untuk semua $t \in 
(a-\epsilon,b+\epsilon)$.  Artinya, buktikan bahwa lintasan mulus dapat diperpanjang
sedikit melampaui ujung-ujung interval parameternya.
\end{exercise}

\begin{exercise} 
Andaikan $\alpha \colon [a,b] \to \R^n$ dan
$\beta \colon [c,d] \to \R^n$ lintasan mulus sepotong-sepotong sedemikian sehingga
$\Gamma \coloneqq \alpha\bigl([a,b]\bigr) = \beta\bigl([c,d]\bigr)$.
Buktikan bahwa terdapat berhingga banyak titik
$\{ p_1,p_2,\ldots,p_k\} \subset \Gamma$ sedemikian sehingga himpunan
$\alpha^{-1}\bigl( \{ p_1,p_2,\ldots,p_k\} \bigr)$
dan
$\beta^{-1}\bigl( \{ p_1,p_2,\ldots,p_k\} \bigr)$
berturut-turut merupakan partisi dari $[a,b]$ dan $[c,d]$ sedemikian sehingga
pembatasan setiap lintasan mulus pada tiap subinterval partisinya sendiri
(artinya, partisi-partisi ini seperti dalam definisi lintasan mulus
sepotong-sepotong).
\end{exercise}

\begin{exercise}
\pagebreak[3]
\leavevmode
\begin{enumerate}[a)]
\item
Andaikan $\gamma \colon [a,b] \to \R^n$ dan $\alpha \colon [c,d] \to \R^n$
dua lintasan mulus yang injektif dan
$\gamma\bigl([a,b]\bigr) = \alpha\bigl([c,d]\bigr)$.  Buktikan bahwa
terdapat reparametrisasi mulus $h \colon [a,b] \to [c,d]$
sedemikian sehingga $\gamma = \alpha \circ h$.\\
Petunjuk 1: Tidak sulit untuk membuktikan bahwa $h$ ada.
Intinya adalah membuktikan bahwa $h$ diferensiabel secara kontinu
dengan turunan tak nol.  Terapkan teorema fungsi implisit,
meskipun pada awalnya dimensinya mungkin tampak tidak sesuai.\\
Petunjuk 2: Tinjau dahulu turunan $h$ pada $(a,b)$.
\item
Buktikan hal yang sama seperti pada bagian (a), tetapi sekarang untuk lintasan tertutup
sederhana, dengan asumsi tambahan $\gamma(a) = \gamma(b) = \alpha(c) = \alpha(d)$.
\item
Buktikan bagian (a) dan (b) untuk lintasan mulus sepotong-sepotong, sehingga diperoleh
reparametrisasi mulus sepotong-sepotong.\\
Petunjuk: Intinya adalah menemukan dua
partisi sedemikian sehingga, pada setiap pasangan subinterval yang bersesuaian,
pembatasan kedua lintasan mempunyai citra yang sama dan keduanya mulus; lihat
latihan sebelumnya.
\end{enumerate}
\end{exercise}

\begin{exercise} 
\pagebreak[1]
Misalkan $\alpha \colon [a,b] \to \R^n$ dan
$\beta \colon [b,c] \to \R^n$ merupakan lintasan mulus sepotong-sepotong dengan
$\alpha(b)=\beta(b)$. Definisikan $\gamma \colon [a,c] \to \R^n$ dengan
\begin{equation*}
\gamma(t) \coloneqq
\begin{cases}
\alpha(t) & \text{jika } t \in [a,b], \\
\beta(t)  & \text{jika } t \in (b,c].
\end{cases}
\end{equation*}
Buktikan bahwa $\gamma$ merupakan lintasan mulus sepotong-sepotong. Buktikan pula
bahwa, jika $\omega$ adalah bentuk-1 yang terdefinisi pada citra lintasan
$\gamma$, maka
\begin{equation*}
\int_{\gamma} \omega =
\int_{\alpha} \omega +
\int_{\beta} \omega .
\end{equation*}
\end{exercise}

\begin{exercise} \label{mv:exercise:closedcurveintegral}
\pagebreak[1]
Misalkan $\gamma \colon [a,b] \to \R^n$ dan
$\beta \colon [c,d] \to \R^n$ merupakan dua lintasan tertutup sederhana yang
mulus sepotong-sepotong. Artinya, $\gamma(a)=\gamma(b)$,
$\beta(c)=\beta(d)$, dan pembatasan $\gamma|_{[a,b)}$ serta
$\beta|_{[c,d)}$ bersifat injektif.
Misalkan $\Gamma = \gamma\bigl([a,b]\bigr) = \beta\bigl([c,d]\bigr)$ dan
$\omega$ merupakan bentuk-1 yang terdefinisi pada $\Gamma \subset \R^n$.
Buktikan bahwa berlaku salah satu dari
\begin{equation*}
\int_\gamma \omega = 
\int_\beta \omega,
\qquad \text{atau} \qquad 
\int_\gamma \omega = 
- \int_\beta \omega.
\end{equation*}
Secara khusus, notasi $\int_{\Gamma} \omega$ bermakna apabila orientasi
penelusuran $\Gamma$ ditentukan.
Petunjuk: Lihat tiga latihan sebelumnya.
\end{exercise}

\begin{exercise} \label{mv:exercise:curveintegral}
Misalkan $\gamma \colon [a,b] \to \R^n$ dan
$\beta \colon [c,d] \to \R^n$ merupakan dua lintasan mulus sepotong-sepotong
yang injektif kecuali pada berhingga banyak titik. Artinya, terdapat himpunan
berhingga $S \subset [a,b]$ dan $T \subset [c,d]$ sedemikian sehingga
$\gamma|_{[a,b]\setminus S}$ dan $\beta|_{[c,d]\setminus T}$ bersifat
injektif. Misalkan
$\Gamma = \gamma\bigl([a,b]\bigr) = \beta\bigl([c,d]\bigr)$ dan $\omega$
merupakan bentuk-1 yang terdefinisi pada $\Gamma \subset \R^n$.

\emph{Koreksi edisi Indonesia.}
Hipotesis di atas saja tidak menjamin
$\int_\gamma\omega=\pm\int_\beta\omega$, berbeda dengan pernyataan dalam teks
sumber. Sebagai contoh, pada kurva berbentuk angka delapan, satu lintasan dapat
membalik orientasi hanya pada salah satu lobus; pilih $\omega$ yang integralnya
tak nol pada kedua lobus. Ambil himpunan berhingga $V\subset\Gamma$ yang memuat
semua titik ujung, titik sudut, dan titik perpotongan-diri kedua parametrisasi,
dan sebut setiap komponen terhubung dari $\Gamma\setminus V$ sebagai busur
reguler. Tambahkan salah satu dari dua hipotesis berikut: kedua lintasan
menginduksi orientasi yang sama pada setiap busur reguler, atau keduanya
menginduksi orientasi yang berlawanan pada setiap busur tersebut. Dalam kedua
kasus itu, buktikan berturut-turut bahwa
\begin{equation*}
\int_\gamma \omega = 
\int_\beta \omega,
\qquad \text{atau} \qquad 
\int_\gamma \omega = 
- \int_\beta \omega.
\end{equation*}
Jika orientasinya bercampur, integral harus diuraikan busur demi busur; sekadar
menentukan satu arah global pada $\Gamma$ tidaklah cukup.
\\
Petunjuk: Terapkan gagasan dari latihan sebelumnya secara terpisah pada setiap
busur reguler.
\end{exercise}

\begin{exercise}
Definisikan $\gamma \colon [0,1] \to \R^2$ dengan
$\gamma(t) \coloneqq \Bigl( t^3 \sin(\nicefrac{1}{t}),\,
t{\bigl(3t^2\sin(\nicefrac{1}{t})-t\cos(\nicefrac{1}{t})\bigr)}^2 \Bigr)$
untuk
$t \neq 0$ dan $\gamma(0) = (0,0)$. Buktikan bahwa
\begin{enumerate}[a)]
\item
$\gamma$ diferensiabel secara kontinu pada $[0,1]$.
\item
Buktikan bahwa himpunan
$\{t\in(0,1] : \gamma'(t)=(0,0)\}$ tak berhingga. Dengan menulis $u=1/t$,
karakterisasikan semua akar melalui persamaan $\tan(u)=u/3$, lalu indekskan
seluruh himpunan tersebut, tanpa ada yang dilewati, sebagai barisan
$t_1>t_2>\cdots>0$ yang konvergen ke~$0$.
\item
Buktikan bahwa titik-titik $\gamma(t_n)$ terletak pada garis $y=0$ dan bahwa
tanda koordinat $x$ dari $\gamma(t_n)$ berselang-seling antara positif dan
negatif.
\item
Buktikan bahwa tidak ada lintasan mulus sepotong-sepotong $\alpha$ yang citranya
sama dengan $\gamma\bigl([0,1]\bigr)$. Petunjuk: Gunakan bagian (c) dan fakta
bahwa $\alpha$ hanya mempunyai berhingga banyak ruas mulus. Pada sedikitnya satu
ruas yang mencapai titik asal, tunjukkan bahwa turunan sepihak harus nol atau
tidak dapat ada; keduanya bertentangan dengan regularitas ruas tersebut.
\item
(Komputer) Jika tersedia perangkat lunak yang dapat menggambar kurva parametrik,
gambarlah kurva ini hanya untuk $t$ dalam rentang $[0,0.1]$; pada
rentang yang lebih lebar perilaku tersebut tidak akan tampak. Secara khusus,
perhatikan bahwa $\gamma\bigl([0,1]\bigr)$ mempunyai tak berhingga banyak
\myquote{sudut}
di dekat titik asal.
\end{enumerate}
Catatan: Anda boleh menggunakan hasil tentang sinus dan kosinus dari kalkulus.
\end{exercise}

%%%%%%%%%%%%%%%%%%%%%%%%%%%%%%%%%%%%%%%%%%%%%%%%%%%%%%%%%%%%%%%%%%%%%%%%%%%%%%

\sectionnewpage
\section{Ketaktergantungan pada lintasan}
\label{sec:pathind}

%mbxINTROSUBSECTION

\sectionnotes{2 kuliah}

\subsection{Integral bebas lintasan}

Misalkan $U \subset \R^n$ himpunan terbuka dan $\omega$ bentuk-1 yang
terdefinisi pada $U$. Integral $\omega$ dikatakan
\emph{\myindex{tidak bergantung pada lintasan}}
apabila, untuk setiap pasangan titik $x,y \in U$ yang dapat dihubungkan oleh
lintasan mulus sepotong-sepotong dan untuk setiap pasangan lintasan semacam itu
$\gamma \colon [a,b] \to U$ dan
$\beta \colon [c,d] \to U$ yang memenuhi $\gamma(a) = \beta(c) = x$
dan $\gamma(b) = \beta(d) = y$, berlaku
\begin{equation*}
\int_\gamma \omega = \int_\beta \omega .
\end{equation*}
Untuk $x$ dan $y$ yang dapat dihubungkan oleh lintasan semacam itu, dalam hal ini
kita cukup menulis
\begin{equation*}
\int_x^y \omega \coloneqq \int_\gamma \omega = \int_\beta \omega .
\end{equation*}
Tidak setiap bentuk-1 menghasilkan integral yang tidak bergantung pada lintasan;
bahkan, kebanyakan tidak demikian.

\begin{example}
Pada $U=\R^2$, tinjau bentuk-1 $\omega=y\,dx$. Misalkan
$\gamma \colon [0,1] \to \R^2$ lintasan
$\gamma(t) \coloneqq (t,0)$ dari $(0,0)$ ke $(1,0)$. Misalkan
$\beta \colon [0,1] \to \R^2$ lintasan
$\beta(t) \coloneqq \bigl(t,(1-t)t\bigr)$ yang juga menghubungkan kedua titik
tersebut. Maka
\begin{align*}
& \int_\gamma y \, dx = 
\int_0^1 \gamma_2(t) \gamma_1'(t) \, dt
=
\int_0^1 0 (1) \, dt = 0 ,\\
& \int_\beta y \, dx = 
\int_0^1 \beta_2(t) \beta_1'(t) \, dt
=
\int_0^1 (1-t)t(1) \, dt = \frac{1}{6} .
\end{align*}
Integral $y\,dx$ bergantung pada lintasan. Secara khusus, notasi
$\int_{(0,0)}^{(1,0)} y\,dx$ tidak mempunyai makna tunggal.
\end{example}

\begin{defn}
Misalkan $U \subset \R^n$ himpunan terbuka dan $f \colon U \to \R$ fungsi
yang diferensiabel secara kontinu. Bentuk-1
\begin{equation*}
df \coloneqq
\frac{\partial f}{\partial x_1} \, dx_1 + 
\frac{\partial f}{\partial x_2} \, dx_2 + \cdots +
\frac{\partial f}{\partial x_n} \, dx_n 
\end{equation*}
disebut \emph{\myindex{turunan total}} dari $f$.

Himpunan terbuka $U \subset \R^n$ disebut
\emph{\myindex{terhubung lintasan}}%
\footnote{Biasanya definisi ini hanya mensyaratkan lintasan kontinu, tetapi
untuk himpunan terbuka kedua definisi tersebut ekuivalen. Lihat latihan.}
apabila, untuk setiap dua titik $x$ dan $y$ di $U$, terdapat lintasan mulus
sepotong-sepotong yang berawal di $x$ dan berakhir di $y$.
\end{defn}

Sebagai latihan, buktikan bahwa setiap himpunan terbuka yang terhubung juga
terhubung lintasan.

\begin{prop} \label{mv:prop:pathinddf}
Misalkan $U \subset \R^n$ himpunan terbuka yang terhubung lintasan dan $\omega$
bentuk-1 yang terdefinisi pada $U$. Integral $\omega$ tidak bergantung pada
lintasan (untuk semua $x,y \in U$) jika dan hanya jika terdapat fungsi
$f \colon U \to \R$ yang diferensiabel secara kontinu sedemikian sehingga
$\omega = df$.

Bahkan, jika fungsi $f$ semacam itu ada, maka untuk setiap pasangan titik
$x,y \in U$ berlaku
\begin{equation*}
\int_{x}^y \omega = f(y)-f(x) .
\end{equation*}
\end{prop}

Dengan kata lain, jika kita menetapkan $p \in U$, maka
$f(x) = C + \int_{p}^x \omega$ untuk suatu konstanta $C$.

\begin{proof}
Pertama, andaikan integral tersebut tidak bergantung pada lintasan. Pilih
$p \in U$. Karena $U$ terhubung lintasan, terdapat lintasan dari $p$ ke setiap
$x \in U$. Definisikan
\begin{equation*}
f(x) \coloneqq \int_{p}^x \omega .
\end{equation*}
Tuliskan $\omega = \omega_1 \,dx_1 + \omega_2 \,dx_2 + \cdots + \omega_n \,dx_n$.
Kita akan membuktikan bahwa, untuk setiap $j = 1,2,\ldots,n$, turunan parsial
$\frac{\partial f}{\partial x_j}$ ada dan sama dengan $\omega_j$.

Ambil sembarang vektor basis standar $e_j$. Karena $U$ terbuka, pilih $r>0$
sedemikian sehingga $\overline{B}(x,r)\subset U$, lalu ambil bilangan real $h$
dengan $0<\sabs{h}<r$. Dengan demikian, seluruh ruas dari $x$ ke $x+he_j$
berada di $U$. Hitunglah
\begin{equation*}
\frac{f(x+h e_j) - f(x)}{h} =
\frac{1}{h} \left( \int_{p}^{x+he_j} \omega - \int_{p}^x \omega \right)
=
\frac{1}{h} \int_{x}^{x+he_j} \omega ,
\end{equation*}
Kesamaan ini mengikuti dari \propref{mv:prop:pathconcat} dan ketaktergantungan
pada lintasan, sebab
$\int_{p}^{x+he_j} \omega =
\int_{p}^{x} \omega +
\int_{x}^{x+he_j} \omega$: kita sengaja memilih lintasan dari $p$ ke
$x+he_j$ yang melalui $x$, lalu memotong lintasan tersebut menjadi dua bagian.
Lihat \figureref{fig:pathindantider}.

\begin{myfigureht}
\myincludepdft{pathindantider}{%
Diagram dengan arah horizontal berlabel e sub j. Sebuah lintasan melengkung
bergerak dari titik p ke titik x, lalu sebuah ruas garis bergerak dari x ke
x ditambah h kali e sub j, sepanjang arah e sub j jika h positif dan arah
sebaliknya jika h negatif.}
\caption{Menggunakan ketaktergantungan pada lintasan untuk menghitung turunan
parsial.\label{fig:pathindantider}}
\end{myfigureht}


Karena integral bebas lintasan, pilih lintasan paling sederhana dari $x$ ke
$x+he_j$, yaitu $\gamma(t) \coloneqq x+t he_j$ untuk $t \in [0,1]$.
Seluruh lintasan ini berada di $U$. Perhatikan bahwa
$\gamma'(t) = h e_j$ hanya mempunyai satu komponen tak nol, yaitu komponen
ke-$j$ yang bernilai $h$. Oleh karena itu,
\begin{equation*}
\frac{1}{h} \int_{x}^{x+he_j} \omega 
=
\frac{1}{h} \int_{\gamma} \omega 
=
\frac{1}{h} \int_0^1 \omega_j(x+the_j) h \, dt 
=
\int_0^1 \omega_j(x+the_j) \, dt  .
\end{equation*}
Kita ingin mengambil limit ketika $h \to 0$. Fungsi $\omega_j$ kontinu di
$x$. Diberikan $\epsilon > 0$, pilih $0<\delta\leq r$ sedemikian sehingga
$\babs{\omega_j(x)-\omega_j(y)} < \epsilon$ setiap kali
$\snorm{x-y}<\delta$. Jika $0<\sabs{h}<\delta$, maka
$\babs{\omega_j(x+the_j)-\omega_j(x)} < \epsilon$ untuk semua
$t \in [0,1]$, dan kita memperoleh taksiran
\begin{equation*}
\abs{\int_0^1 \omega_j(x+the_j) \, dt  - \omega_j(x)}
=
\abs{\int_0^1 \bigl( \omega_j(x+the_j) - \omega_j(x) \bigr) \, dt}
\leq
\epsilon .
\end{equation*}
Jadi,
\begin{equation*}
\lim_{h\to 0}\frac{f(x+h e_j) - f(x)}{h} = \omega_j(x) .
\end{equation*}
Semua turunan parsial $f$ ada, dan setiap
$\frac{\partial f}{\partial x_j}$ sama dengan koefisien yang bersesuaian,
$\omega_j$, yang kontinu. Jadi, $f$ diferensiabel secara kontinu dan
$df = \omega$.

Untuk arah sebaliknya, andaikan terdapat fungsi $f$ yang diferensiabel secara
kontinu dan memenuhi $df = \omega$. Ambil lintasan mulus
$\gamma \colon [a,b] \to U$ dengan $\gamma(a) = x$ dan
$\gamma(b) = y$. Maka
\begin{equation*}
\begin{split}
\int_\gamma df
& =
\int_a^b
\biggl(
\frac{\partial f}{\partial x_1}\bigl(\gamma(t)\bigr) \gamma_1'(t)+
\frac{\partial f}{\partial x_2}\bigl(\gamma(t)\bigr) \gamma_2'(t)+ \cdots +
\frac{\partial f}{\partial x_n}\bigl(\gamma(t)\bigr) \gamma_n'(t)
\biggr) \, dt
\\
& = 
\int_a^b
\frac{d}{dt} \Bigl[ f\bigl(\gamma(t)\bigr) \Bigr]\, dt
\\
& = f(y)-f(x) .
\end{split}
\end{equation*}
Nilai integral hanya bergantung pada $x$ dan $y$, bukan pada lintasan yang
ditempuh. Untuk lintasan mulus sepotong-sepotong dengan partisi
$a=t_0<t_1<\cdots<t_k=b$, terapkan perhitungan di atas pada setiap pembatasan
$\gamma|_{[t_{i-1},t_i]}$, lalu jumlahkan. Suku-suku pada titik partisi saling
meniadakan, sehingga tetap diperoleh
$\int_\gamma df=f\bigl(\gamma(b)\bigr)-f\bigl(\gamma(a)\bigr)$. Jadi, integral
tersebut bebas lintasan.
\end{proof}

Ketaktergantungan pada lintasan dapat dinyatakan dengan lebih ringkas melalui
integral pada lintasan tertutup.

\begin{prop}
Misalkan $U \subset \R^n$ himpunan terbuka yang terhubung lintasan dan $\omega$
bentuk-1 yang terdefinisi pada $U$. Maka $\omega = df$ untuk suatu fungsi
$f \colon U \to \R$ yang diferensiabel secara kontinu jika dan hanya jika
\begin{equation*}
\int_{\gamma} \omega = 0
\qquad
\text{untuk setiap lintasan tertutup mulus sepotong-sepotong }
\gamma \colon [a,b] \to U.
\end{equation*}
\end{prop}

\begin{proof}
Andaikan $\omega = df$ dan misalkan $\gamma$ lintasan tertutup yang mulus
sepotong-sepotong. Karena $\gamma(a) = \gamma(b)$, proposisi sebelumnya
memberikan
\begin{equation*}
\int_{\gamma} \omega = f\bigl(\gamma(b)\bigr) - f\bigl(\gamma(a)\bigr) = 0 .
\end{equation*}

Sekarang andaikan $\int_{\gamma} \omega = 0$ untuk setiap lintasan tertutup
mulus sepotong-sepotong $\gamma$. Ambil dua titik $x,y$ di $U$ dan dua lintasan
mulus sepotong-sepotong yang menghubungkannya. Setelah reparametrisasi afin yang
mempertahankan orientasi, tanpa mengurangi keumuman tuliskan kedua lintasan itu
sebagai $\alpha \colon [0,1] \to U$ dan $\beta \colon [0,1] \to U$, dengan
$\alpha(0) = \beta(0) = x$ dan $\alpha(1) = \beta(1) = y$. Lihat
\figureref{fig:twopaths}.
\begin{myfigureht}
\myincludepdft{twopaths}{%
Dua titik, x dan y, beserta dua lintasan berbeda, alfa dan beta, yang sama-sama
berawal di x dan berakhir di y.}
\caption{Dua lintasan dari $x$ ke $y$.\label{fig:twopaths}}
\end{myfigureht}

Definisikan $\gamma \colon [0,2] \to U$ dengan
\begin{equation*}
\gamma(t) \coloneqq
\begin{cases}
\alpha(t)  & \text{jika } t \in [0,1], \\
\beta(2-t) & \text{jika } t \in (1,2].
\end{cases}
\end{equation*}
Lintasan ini mulus sepotong-sepotong: gabungkan partisi berhingga untuk
$\alpha$ dengan partisi lintasan balik $t\mapsto\beta(2-t)$. Secara khusus,
$\gamma|_{[0,1]}(t) = \alpha(t)$ dan
$\gamma|_{[1,2]}(t) = \beta(2-t)$, serta
$\gamma(1) = \alpha(1) = \beta(2-1)$. Lintasan ini juga tertutup karena
$\gamma(0) = \alpha(0) = \beta(0) = \gamma(2)$. Jadi,
\begin{equation*}
0 = \int_{\gamma} \omega = \int_{\alpha} \omega - \int_{\beta} \omega .
\end{equation*}
Kesamaan ini pertama-tama mengikuti dari \propref{mv:prop:pathconcat}; bagian
kedua menelusuri $\beta$ dalam arah sebaliknya sehingga memberikan
$-\int_\beta\omega$. Dengan demikian, integral $\omega$ pada $U$ bebas
lintasan.
\end{proof}

Apa pun perumusan ketaktergantungan pada lintasan, kriterianya sering sulit
diperiksa---kita harus memeriksa sesuatu \myquote{untuk semua lintasan.}
Ada kriteria lokal yang lebih mudah, berupa persamaan diferensial, yang menjamin
ketaktergantungan pada lintasan. Dengan kata lain, kriteria itu menjamin adanya
\emph{\myindex{antiturunan}} $f$ yang turunan totalnya adalah bentuk-1
$\omega$ yang diberikan. Karena kriterianya lokal, pada umumnya kita hanya
memperoleh fungsi $f$ secara lokal.
Jika kriteria lokal tersebut terpenuhi, antiturunan dapat ditemukan pada setiap
domain
\emph{\myindex{terhubung sederhana}}. Secara informal, ini adalah himpunan
terbuka terhubung tempat setiap dua lintasan dengan titik-titik ujung yang sama
dapat dihomotopikan satu sama lain di dalam himpunan tersebut sambil
mempertahankan kedua titik ujung. Untuk menyederhanakan pembahasan, kita
membuktikan hasil ini bagi domain berbentuk bintang, yang sering kali sudah
memadai. Sebagai tambahan, pembuktian untuk kasus berbentuk bintang membangun
antiturunannya secara eksplisit. Karena bola berbentuk bintang, hasil ini pun
berlaku secara lokal.

\begin{defn}
Misalkan $U \subset \R^n$ himpunan terbuka dan $p \in U$. Kita menyebut $U$
\emph{\myindex{domain berbentuk bintang}}
terhadap $p$ apabila, untuk setiap titik lain $x \in U$, ruas garis $[p,x]$
berada di dalam $U$, yaitu apabila
$(1-t)p + tx \in U$ untuk semua $t \in [0,1]$.
Jika kita hanya mengatakan \emph{berbentuk bintang}, yang dimaksud adalah bahwa
$U$ berbentuk bintang terhadap suatu $p \in U$. Lihat
\figureref{fig:starshaped}.
\begin{myfigureht}
\myincludepdft{starshaped}{%
Daerah berarsir berbentuk bintang terhadap titik p. Beberapa titik di dalam
daerah, termasuk titik x, dihubungkan ke p, dan setiap ruas garis tersebut
seluruhnya berada di dalam daerah.}
\caption{Domain berbentuk bintang terhadap $p$.\label{fig:starshaped}}
\end{myfigureht}
\end{defn}

Perhatikan perbedaan antara berbentuk bintang dan konveks. Setiap himpunan
konveks berbentuk bintang, tetapi domain berbentuk bintang belum tentu konveks.

\begin{thm}[\myindex{Lema Poincar\'e}]
Misalkan $U \subset \R^n$ domain berbentuk bintang dan $\omega$ bentuk-1 yang
diferensiabel secara kontinu serta terdefinisi pada $U$. Artinya,
\begin{equation*}
\omega =
\omega_1 \,dx_1 +
\omega_2 \,dx_2 + \cdots +
\omega_n \,dx_n ,
\end{equation*}
dengan $\omega_1,\omega_2,\ldots,\omega_n$ fungsi-fungsi yang diferensiabel
secara kontinu. Andaikan, untuk setiap $j$ dan $k$,
\begin{equation*}
\frac{\partial \omega_j}{\partial x_k} = \frac{\partial \omega_k}{\partial x_j} .
\avoidbreak
\end{equation*}
Maka terdapat fungsi $f \colon U \to \R$ yang diferensiabel secara kontinu
hingga orde dua dan memenuhi $df = \omega$.
\end{thm}

Syarat pada turunan $\omega$ tepat menyatakan bahwa urutan diferensiasi
parsial kedua dapat dipertukarkan. Jika $df = \omega$ dan $f$ diferensiabel
secara kontinu hingga orde dua, maka
\begin{equation*}
\frac{\partial \omega_j}{\partial x_k}
=
\frac{\partial^2 f}{\partial x_k \partial x_j} 
=
\frac{\partial^2 f}{\partial x_j \partial x_k} 
=
\frac{\partial \omega_k}{\partial x_j} .
\end{equation*}
Syarat ini jelas perlu. Lema Poincar\'e menyatakan bahwa syarat tersebut juga
cukup apabila $U$ berbentuk bintang.

\begin{proof}
Andaikan $U$ berbentuk bintang terhadap
$p=(p_1,p_2,\ldots,p_n) \in U$. Untuk
$x = (x_1,x_2,\ldots,x_n) \in U$, definisikan pemetaan radial
$\gamma_x \colon [0,1] \to U$ dengan
$\gamma_x(t) \coloneqq (1-t)p + tx$, sehingga $\gamma_x'(t) = x-p$. Jika
$x\neq p$, pemetaan ini merupakan lintasan mulus; jika $x=p$, gunakan
perluasan definisi integral bagi pemetaan $C^1$, yang memberi integral nol.
Definisikan
\begin{equation*}
f(x) \coloneqq \int_{\gamma_x} \omega
=
\int_0^1
\left(
\sum_{k=1}^n
\omega_k \bigl((1-t)p + tx \bigr) \, (x_k-p_k)
\right) dt .
\end{equation*}
Untuk membenarkan diferensiasi di bawah tanda integral, tetapkan
$x\in U$ dan pilih $\rho>0$ sedemikian sehingga
$\overline{B}(x,\rho)\subset U$. Karena $U$ berbentuk bintang terhadap $p$,
himpunan kompak
$\{(1-t)p+ty:t\in[0,1],\ y\in\overline{B}(x,\rho)\}$ seluruhnya berada di
$U$. Koefisien $\omega_k$ beserta semua turunan parsial pertamanya kontinu pada
himpunan tersebut. Karena itu, kita boleh mendiferensialkan terhadap $x_j$ di
bawah tanda integral:
\begin{equation*}
\begin{split}
\frac{\partial f}{\partial x_j}(x) & =
\int_0^1
\left(
\left(
\sum_{k=1}^n
\frac{\partial \omega_k}{\partial x_j} \bigl((1-t)p + tx \bigr) \, t
(x_k-p_k)
\right)
+
\omega_j \bigl((1-t)p + tx \bigr)
\right)
 dt
\\
& = 
\int_0^1
\left(
\left(
\sum_{k=1}^n
\frac{\partial \omega_j}{\partial x_k} \bigl((1-t)p + tx \bigr) \, t
(x_k-p_k)
\right)
+
\omega_j \bigl((1-t)p + tx \bigr)
\right) dt
\\
& = 
\int_0^1
\frac{d}{dt}
\Bigl[
t \omega_j\bigl((1-t)p + tx \bigr)
\Bigr]
\,
dt
\\
&= \omega_j(x) .
\end{split}
\end{equation*}
Inilah kesamaan yang hendak dibuktikan. Karena
$\frac{\partial f}{\partial x_j}=\omega_j$ dan setiap $\omega_j$ diferensiabel
secara kontinu, semua turunan parsial kedua $f$ ada dan kontinu. Jadi, $f$
diferensiabel secara kontinu hingga orde dua dan $df=\omega$.
\end{proof}

\begin{example}
Tanpa hipotesis tertentu pada $U$, teorema tersebut tidak benar. Pada
$\R^2 \setminus \{ (0,0) \}$, definisikan
\begin{equation*}
\omega(x,y) \coloneqq \frac{-y}{x^2+y^2} \,dx + \frac{x}{x^2+y^2} \,dy .
\end{equation*}
Maka
\begin{equation*}
\frac{\partial}{\partial y} \left[ 
\frac{-y}{x^2+y^2} \right] =
\frac{y^2-x^2}{{(x^2+y^2)}^2}
=
\frac{\partial}{\partial x} \left[ 
\frac{x}{x^2+y^2} \right] .
\end{equation*}
Namun, tidak ada fungsi
$f \colon \R^2 \setminus \{ (0,0) \} \to \R$ yang memenuhi $df = \omega$.
Dalam \exampleref{example:mv:irrotoneformint}, kita mengintegralkan dari
$(1,0)$ kembali ke $(1,0)$ sepanjang lingkaran satuan berlawanan arah jarum
jam, yaitu dengan
$\gamma(t) = \bigl(\cos(t),\sin(t)\bigr)$ untuk $t \in [0,2\pi]$, dan
memperoleh integral $2\pi$. Jika integral itu bebas lintasan---atau, dengan kata
lain, jika terdapat $f$ yang memenuhi $df = \omega$---hasilnya seharusnya nol.
\end{example}

\subsection{Medan vektor}

Jika pada setiap titik terdapat sebuah vektor---yang kita sebut medan
vektor---kita dapat mengintegralkan komponennya sepanjang suatu lintasan.
Integral semacam ini, misalnya, menghitung usaha yang dilakukan medan gaya pada
partikel yang bergerak sepanjang lintasan.

\begin{defn}
Misalkan $U \subset \R^n$ suatu himpunan. Fungsi kontinu
$v \colon U \to \R^n$ disebut \emph{\myindex{medan vektor}}. Tuliskan
$v = (v_1,v_2,\ldots,v_n)$.

Diberikan lintasan mulus $\gamma \colon [a,b] \to \R^n$ dengan
$\gamma\bigl([a,b]\bigr) \subset U$, kita mendefinisikan integral lintasan
medan vektor $v$ sebagai
\glsadd{not:pathintegralvecfield}
\begin{equation*}
\int_{\gamma} v \cdot d\gamma
\coloneqq
\int_a^b v\bigl(\gamma(t)\bigr) \cdot \gamma'(t) \, dt ,
\end{equation*}
di sini tanda titik menyatakan hasil kali titik standar. Untuk lintasan mulus
sepotong-sepotong, definisinya diperoleh dengan mengintegralkan pada setiap
interval mulus lalu menjumlahkan hasil-hasilnya.
\end{defn}

Dengan menguraikan definisi, kita memperoleh
\begin{equation*}
\int_{\gamma} v \cdot d\gamma
=
\int_{\gamma} \bigl(v_1 \,dx_1 + v_2 \,dx_2 + \cdots + v_n \,dx_n\bigr) .
\end{equation*}
Semua yang kita ketahui tentang integrasi bentuk-1 berlaku pula untuk integrasi
medan vektor. Integral medan vektor disebut bebas lintasan apabila, untuk setiap
dua lintasan mulus sepotong-sepotong $\gamma$ dan $\beta$ di $U$ dengan titik
awal dan titik akhir yang sama, berlaku
\begin{equation*}
\int_\gamma v \cdot d\gamma = \int_\beta v \cdot d\beta .
\end{equation*}
Jika $U$ terbuka dan terhubung lintasan, sifat ini berlaku jika dan hanya jika
terdapat fungsi
$f \colon U\to\R$ yang diferensiabel secara kontinu dengan
$v = \nabla f$. Artinya, $v$ merupakan gradien suatu fungsi. Fungsi $f$
kemudian disebut \emph{\myindex{potensial}} bagi $v$.

Medan vektor $v$ yang integral lintasannya bebas lintasan disebut
\emph{\myindex{medan vektor konservatif}}. Istilah ini digunakan karena medan
vektor semacam itu muncul dalam sistem fisika tempat suatu besaran, yaitu
energi, bersifat kekal.

\subsection{Latihan}

\begin{exercise}
Carilah $f \colon \R^2 \to \R$ sedemikian sehingga $df = xe^{x^2+y^2}\, dx +
ye^{x^2+y^2} \, dy$.
\end{exercise}

\begin{exercise}
Carilah $\omega_2 \colon \R^2 \to \R$ sedemikian sehingga terdapat fungsi
$f \colon \R^2 \to \R$ yang diferensiabel secara kontinu dan memenuhi
$df = e^{xy} \,dx + \omega_2 \,dy$.
\end{exercise}

\begin{exercise}
Jabarkan secara mandiri rincian kasus lintasan mulus sepotong-sepotong dalam
pembuktian \propref{mv:prop:pathinddf}: terapkan perhitungan bagi lintasan mulus
pada setiap subinterval partisi, lalu tunjukkan bahwa suku-suku di titik partisi
saling meniadakan.
\end{exercise}

\begin{exercise}
Buktikan bahwa jika $U \subset \R^n$ merupakan domain berbentuk bintang, maka
$U$ terhubung lintasan. Petunjuk: Untuk dua titik berbeda, gabungkan ruas-ruas
garis melalui pusat bintang dan hilangkan setiap ruas konstan. Untuk
menghubungkan suatu titik $x$ kembali ke dirinya sendiri, pilih $q\neq x$ dalam
bola kecil yang termuat di $U$, lalu tempuh ruas dari $x$ ke $q$ dan kembali ke
$x$.
\end{exercise}

\begin{exercise}
Buktikan bahwa
$U \coloneqq \R^2 \setminus \{ (x,y) \in \R^2 : x \leq 0, y=0 \}$ merupakan
domain berbentuk bintang, lalu carilah semua titik $(x_0,y_0) \in U$ yang
menjadikan $U$ berbentuk bintang terhadap $(x_0,y_0)$.
\end{exercise}

\begin{exercise}
Misalkan $U_1$ dan $U_2$ dua himpunan terbuka di $\R^n$ dengan
$U_1 \cap U_2$ tak kosong dan terhubung lintasan. Andaikan terdapat fungsi
$f_1 \colon U_1 \to \R$ dan $f_2 \colon U_2 \to \R$ yang keduanya
diferensiabel secara kontinu hingga orde dua serta memenuhi
$df_1 = df_2$ pada $U_1 \cap U_2$. Buktikan bahwa terdapat fungsi
$F \colon U_1 \cup U_2 \to \R$ yang diferensiabel secara kontinu hingga orde
dua sedemikian sehingga $dF = df_1$ pada $U_1$ dan $dF = df_2$ pada $U_2$.
\end{exercise}

\begin{exercise}[Sulit]
Misalkan $\gamma \colon [a,b] \to \R^n$ lintasan sederhana, tidak tertutup,
dan mulus sepotong-sepotong (jadi, $\gamma$ injektif). Andaikan $\omega$
bentuk-1 yang diferensiabel secara kontinu dan terdefinisi pada suatu himpunan
terbuka $V$, dengan $\gamma\bigl([a,b]\bigr) \subset V$ dan
$\frac{\partial \omega_j}{\partial x_k} = \frac{\partial \omega_k}{\partial
 x_j}$
untuk semua $j$ dan $k$. Buktikan bahwa terdapat himpunan terbuka $U$ dengan
$\gamma\bigl([a,b]\bigr) \subset U \subset V$ dan fungsi
$f \colon U \to \R$ yang diferensiabel secara kontinu hingga orde dua serta
memenuhi $df = \omega$.
\\
Petunjuk 1: $\gamma\bigl([a,b]\bigr)$ kompak.\\
Petunjuk 2: Tutupi citra lintasan, menurut urutan parameternya, dengan berhingga
banyak bola terbuka $B_1,\ldots,B_m\subset V$ sedemikian sehingga
$B_i\cap B_{i+1}\neq\varnothing$ dan $B_i\cap B_j=\varnothing$ jika
$\sabs{i-j}>1$.\\
Petunjuk 3: Setiap $B_i\cap B_{i+1}$ konveks, sehingga terhubung lintasan.
Gunakan latihan sebelumnya secara induktif untuk merekatkan antiturunan lokal
pada bola-bola tersebut.
\end{exercise}

\begin{exercise}
\pagebreak[2]
\leavevmode
\begin{enumerate}[a)]
\item
Buktikan bahwa himpunan terbuka yang terhubung $U \subset \R^n$ juga terhubung
lintasan. Petunjuk: Mulailah dengan titik $x \in U$, lalu definisikan
$U_x \subset U$ sebagai himpunan semua titik yang dapat dicapai dari $x$ oleh
suatu lintasan. Buktikan bahwa $U_x$ dan $U \setminus U_x$ keduanya terbuka.
Untuk melihat bahwa $x\in U_x$, pilih $q\neq x$ dalam bola kecil yang termuat di
$U$ dan gunakan lintasan pergi-pulang sepanjang ruas $xq$. Jadi, $U_x$ tak
kosong dan haruslah $U_x = U$.
\item
Buktikan arah sebaliknya, yaitu bahwa setiap himpunan
$U \subset \R^n$ yang terhubung lintasan juga terhubung.\footnote{Koreksi edisi
Indonesia: Teks sumber menambahkan syarat \myquote{terbuka}, tetapi arah ini
tidak memerlukannya karena setiap lintasan mulus sepotong-sepotong bersifat
kontinu. Keterbukaan diperlukan pada bagian (a) dan pada kesetaraan dengan
definisi berbasis lintasan kontinu dalam latihan berikutnya.} Petunjuk: Untuk
memperoleh kontradiksi, andaikan terdapat dua himpunan terbuka relatif yang tak
kosong dan saling lepas dengan gabungan $U$. Lalu gunakan lintasan mulus
sepotong-sepotong (dan karenanya kontinu) dari suatu titik di himpunan pertama
ke suatu titik di himpunan kedua.
\end{enumerate}
\end{exercise}

\begin{exercise}
Biasanya keterhubungan lintasan didefinisikan dengan lintasan kontinu, bukan
lintasan mulus sepotong-sepotong. Buktikan bahwa kedua definisi itu ekuivalen
untuk himpunan terbuka di $\R^n$. Dengan kata lain, buktikan pernyataan berikut.\\
Andaikan $U \subset \R^n$ terbuka dan untuk setiap $x,y \in U$ terdapat fungsi
kontinu $\gamma \colon [a,b] \to U$ dengan $\gamma(a) = x$ dan
$\gamma(b) = y$. Maka $U$ terhubung lintasan menurut definisi kita, yaitu
terdapat lintasan mulus sepotong-sepotong di $U$ dari $x$ ke $y$. Untuk
$x\neq y$, gunakan kekompakan citra lintasan kontinu untuk membangun rantai
poligonal berhingga di dalam bola-bola yang termuat di $U$, lalu hapus
pengulangan titik sudut yang berurutan. Untuk $x=y$, gunakan lintasan pergi-pulang di
dalam bola kecil di $U$.
\end{exercise}

\begin{exercise}[Sulit]
\pagebreak[2]
Tinjau
\begin{equation*}
\omega(x,y) \coloneqq \frac{-y}{x^2+y^2} \, dx + \frac{x}{x^2+y^2} \, dy .
\end{equation*}
Bentuk-1 ini terdefinisi pada $\R^2 \setminus \{ (0,0) \}$. Misalkan
$\gamma \colon [a,b] \to \R^2 \setminus \{ (0,0) \}$ lintasan tertutup yang
mulus sepotong-sepotong. Definisikan
$R \coloneqq \{ (x,y) \in \R^2 : x \leq 0 \text{ dan } y=0 \}$, yaitu sinar
sumbu-$x$ nonpositif.

\emph{Koreksi edisi Indonesia.} Andaikan himpunan parameter
$T\coloneqq\{t\in[a,b):\gamma(t)\in R\}$, bukan sekadar himpunan titik citra
pada $R$, berhingga dan mempunyai tepat $k$ anggota. Buktikan bahwa
\begin{equation*}
\int_{\gamma} \omega = 2 \pi \ell ,
\end{equation*}
untuk suatu bilangan bulat $\ell$ dengan $\sabs{\ell} \leq k$.\\
Petunjuk 1: Buktikan terlebih dahulu bahwa, untuk lintasan
$\beta\colon[c,d]\to\R^2\setminus\{(0,0)\}$ dengan
$\beta(c),\beta(d)\in R$ dan $\beta((c,d))\cap R=\varnothing$, integral
$\int_{\beta} \omega$ bernilai $-2\pi$, 0, atau $2\pi$.
\\
Petunjuk 2: Di atas Anda telah membuktikan bahwa $\R^2 \setminus R$ berbentuk
bintang.
\\
Catatan: Bilangan $\ell$ disebut \emph{\myindex{bilangan putar}}; bilangan ini
menghitung banyaknya putaran bersih $\gamma$ mengelilingi titik asal; putaran
berlawanan arah jarum jam dihitung positif.
\end{exercise}
